\documentclass[11pt]{article}
\usepackage[utf8]{inputenc}
\usepackage[margin=1in]{geometry}
\usepackage{graphicx} % Required for inserting images
\usepackage{amsfonts}
\usepackage{amsmath}
\usepackage{dirtytalk}
\usepackage{algorithm}
\usepackage{mathtools}
\usepackage{xcolor}
\usepackage{amssymb}
\usepackage[title]{appendix}
\usepackage[noend]{algpseudocode}
\usepackage{amsthm}
\usepackage{bigints}
\usepackage{thm-restate}
\usepackage[]{hyperref}
\usepackage{cleveref}
\usepackage{appendix}
\usepackage{makecell}
\usepackage{bigints}
\usepackage{mathtools}
\usepackage{subfigure}
\usepackage{booktabs}
\usepackage{bbding}
\usepackage{wrapfig}
\usepackage{bbm}
\usepackage{authblk}
\usepackage{cleveref}

\newtheorem{thm}{Theorem}[section]

% everyone else shares the same counter
\newtheorem{lemma}[thm]{Lemma}
\newtheorem{prop}[thm]{Proposition}
\newtheorem{corollary}[thm]{Corollary}
\newtheorem{asmpt}[thm]{Assumption}
\newtheorem{defn}[thm]{Definition}
\newtheorem{rmk}[thm]{Remark}
\newtheorem{claim}[thm]{Claim}

\usepackage[numbers]{natbib}
\bibliographystyle{unsrtnat}


\newcommand{\expnumber}[2]{{#1}\mathrm{e}{#2}}
\newcommand{\man}{\mathcal{M}}
\newcommand{\norm}[1]{\left\lVert#1\right\rVert}


\title{Locally Linear Single Index Regression}

\author{Tristan Luca Saidi}
\author{Arun Kuchibhotla}

\affil[1]{Department of Statistics and Data Science, Carnegie Mellon University}
\begin{document}


\maketitle

\textcolor{red}{Saving this version because in the next version I'm gonna get rid of the requirement that $p^*$ is supported on $\mathcal{B}(0, \epsilon)$. I had this originally because I thought it was required to invoke the Newey and Ruud lemma, but since we're not using that directly anymore I believe its okay to get rid of it.}

\section{Introduction}
\subsection{Single-Index Setting}
We consider the regression setting where we are given samples $\{(x_1, y_1), \dots, (x_d, y_n)\}$ from some joint distribution over $\mathbb{R}^d\times \mathbb{R}$ with bounded support. In the single-index setting, we know that the conditional expectation of $y$ given $x$ has the following compositional structure,
\[y_i = f(\Pi(x_i)) + \varepsilon_i \qquad \text{with } \qquad \mathbb{E}[\varepsilon_i] = 0, \text{Var}(\varepsilon_i) = \sigma^2.\]
where $\Pi:\mathbb{R}^d \rightarrow \Gamma$ and $f: \Gamma\rightarrow\mathbb{R}$. The classical single-index setting assumes that $\Gamma = \mathbb{R}$, $\Pi(x) = b^{T}x$ for some $b \in \mathbb{R}^d$, and thus can be thought of as a projection onto a line. In this work, we consider the setting where $\Pi$ is a closest-point projection onto a compact regular $C^3$ curve $\gamma$. We choose to represent $\gamma$ as a function $\gamma:I \rightarrow \Gamma \subset \mathbb{R}^d$ with $I = [0,\text{len}_{\gamma}]$, where $\gamma$ is parameterized by arclength $t$. We assume that $\gamma$ is injective, which ensures that $\gamma^{-1}:\Gamma\rightarrow I$ is well defined and rules out any self-intersections. We define the closest-point projection map as
\begin{equation}\label{def: projection}
\Pi(x) = \arg\inf_{t\in I}\|x - \gamma(t)\|_2.
\end{equation}
\begin{restatable}[Infimum is attained]{rmk}{} \label{remark: infimum}
    The infimum in \cref{def: projection} is attained by at least one point in the set $I$ for every $x \in \mathbb{R}^d$.
\end{restatable}
\noindent \textit{Proof.} Since the set $\Gamma$ is compact and the Euclidean distance $\|x-a\|_2$ is continuous, the infimum must be attained on some element in $\Gamma$. The result follows from the injectivity of $\gamma$. 

\qed{}

\subsection{Geometry of curves in $\mathbb{R}^d$}
Before presenting our estimator and its properties, we will remark on relevant results from differential geometry and geometric measure theory. 


\begin{restatable}[Medial axis, \cite{federer1959curvature}]{defn}{}
    Given a closed subset $A \subset \mathbb{R}^d$, the medial axis $M(A)$ is the subset of $\mathbb{R}^d$ consisting of points with at least two nearest neighbors on $A$. Formally,
    \begin{equation}
    M(A) = \{x \in \mathbb{R}^d : |\{y \in A : \|x-y\|_2 = d(x,A)\}| > 1\}. \label{eq: medial axis definition}
    \end{equation}
\end{restatable}

\begin{restatable}[Reach, \cite{federer1959curvature}]{defn}{}
    The reach of a closed subset $A \subset \mathbb{R}^d$ is defined as 
    \begin{equation}
        \tau_A = \inf_{p\in A} d(p,M(A)) = \inf_{z\in M(A)}d(z,A). \label{eq: reach definition}
    \end{equation}
\end{restatable}


\begin{restatable}[Boundedness of $\gamma^{(k)}$]{rmk}{} \label{remark: deriv bound}
    The curve $\gamma$ has bounded first, second and third derivatives. Namely, for any $\gamma \in C^3(I)$ and $k \in \{1,2,3\}$ we have 
    \[\sup_{t\in I}\|\gamma^{(k)}(t)\|_2 \leq M_k\]
    for some $|M_k| < \infty$.
\end{restatable}
\noindent \textit{Proof.} Since $\gamma \in C^3(I)$, $\gamma^{(k)}$ is continuous for $k \leq 3$. Since $I$ is a compact set, by the Extreme Value Theorem $\gamma^{(k)}$ must be bounded on $I$. 

\qed{}


\noindent We will now discuss some properties of the closest point projection map $\Pi$. 


\begin{restatable}[Differentiability of $\Pi$]{thm}{} \label{thm: differentiability of pi}
    There exists an \textit{open} set $R$ of full $d$-dimensional Lebesgue measure such that the following conditions hold for all $x \in R$:
    \begin{enumerate}
        \item $s = \Pi(x)$ is unique and well defined
        \item both $\nabla\Pi(x)$, $\nabla^2\Pi(x)$ exist and can be expressed as $\nabla\Pi(x)  = -\gamma'(s)/g(x)$ and 
        \begin{equation}
            \nabla^2\Pi(x) = \frac{1}{g^2(x)}\left(\gamma''(s)\gamma'(s)^T + \gamma'(s)\gamma''(s)^T\right)- \frac{\langle\gamma^{(3)}(s), x - \gamma(s)\rangle}{g^3(x)}\gamma'(s)\gamma'(s)^T
        \end{equation}
        where $g(x) = -1 + \langle\gamma''(s), x-\gamma(s)\rangle$. 
    \end{enumerate}
\end{restatable}
\noindent \textit{Proof.} Let $B$ and $\overline{\text{Sing}(\Gamma)}$ be defined as in \cref{eq: ift degeneracy set} and \cref{eq: singular set} respectively. By \Cref{lemma: IFT degen} and \Cref{lemma: singular set measure} we know that both sets have Lebesgue measure $0$. By \Cref{lemma: IFT degen closed} we also know that $B$ is closed. It follows that $B \, \cup \overline{\text{Sing}(\Gamma)}$ has measure $0$ and is also closed, from which we have that 
\begin{equation}
    R = \left(B \, \cup \overline{\text{Sing}(\Gamma)}\right)^c \label{eq: good set}
\end{equation}
is open with full measure. We also know that for any $x \in R$
\begin{enumerate}
    \item $x \notin \overline{\text{Sing}(\Gamma)}$, which means $\Pi(x)$ is differentiable at $x$ and $\Pi(x)$ therefore exists and is unique. 
    \item $x \notin B$, which allows us to apply \Cref{prop: proj gradient} and \Cref{prop: proj hessian} to get the expressions for $\nabla\Pi(x)$ and $\nabla^2\Pi(x)$ in the theorem statement.
\end{enumerate}
\qed{}

% \begin{restatable}[Lebesgue almost everywhere uniqueness of $\Pi$, \cite{erdos1945some}]{corollary}{} \label{remark: uniqueness}
%     The closest point projection $\Pi$ onto the closed regular $C^3$ curve $\gamma$ is Lebesgue almost everywhere unique. 
% \end{restatable}


\subsection{Distribution of Covariates}
Previous works utilize the Linear Conditional Expectation (LCE) property, where we assume that the following condition on the distribution of the covariates holds: for any $b \in \mathbb{R}^p$ and $B \in \mathbb{R}^{d\times p}$, then there exists a $c \in \mathbb{R}^d$ such that the conditional expectation $\mathbb{E}[b^{T}x \, | B^{T}x] = B^{T}xc$. In this work we will only assume that the support of $x$ is contained in $\mathcal{B}(0, r)$ for some $r > 0$. We will also suppose that $\Gamma  = \text{Im}(\gamma) \subset \mathcal{B}(0, r)$ as well. 

\begin{restatable}[Elliptically symmetric distributions]{defn}{}\label{def: elliptical symmetry}
Let $\mu \in \mathbb{R}^d$ and let $\Lambda \in \mathbb{R}^{d\times d}$ with $\Lambda \succeq 0$. If $x$ has density
\begin{equation}
     f(x) = |\Lambda|^{-1/2}g\left((x-m)^{T}\Lambda^{-1}(x - m)\right)
\end{equation}
and $g:[0, \infty) \rightarrow[0, \infty)$ does not depend on $m$, $\Lambda$, then the distribution of $x$ is said be elliptically symmetric. 
\end{restatable}

\begin{restatable}[Spherically symmetric distributions]{defn}{}\label{def: spherical symmetry}
If $X\in \mathbb{R}^d$ has density satisfying \Cref{def: elliptical symmetry} with $\Lambda = I_d$, then the distribution of $x$ is said to be spherically symmetric. Equivalently, a random variable $x$ has a spherically symmetric distribution about its mean $m$ if for every orthogonal matrix $Q$, we have $Q(x - m) \overset{d}{=} x - m$ (where $\overset{d}{=}$ denotes equivalence in distribution). 
\end{restatable}

\subsection{Recovering the underlying curve}
For a given $x_0 \in \mathbb{R}^d$ will construct an estimator $\widehat{\beta}_{x_0} \approx  \begin{bmatrix} \alpha_0 & \alpha_1\gamma'(\Pi(x_0)) \end{bmatrix}^T$ for some $\alpha_0, \alpha_1 \in \mathbb{R}$ using the following procedure. Let $\epsilon > 0$ and define
\begin{align} 
\widehat{\beta}_{x_0,\epsilon}, \widehat{\alpha}_{x_0, \epsilon} &= \arg\min_{\beta \in \mathbb{R}^d,\, \alpha \in \mathbb{R}} \sum_{x_i \in \mathcal{B}(x_0, \epsilon)}\frac{p^*(x_i - x_0;\theta_0, \epsilon)}{\hat{p}(x_i)}\left(y_i - (x_i - x_0)^{T}\beta - \alpha\right)^2 \\
&= \arg\min_{\beta \in \mathbb{R}^d,\, \alpha \in \mathbb{R}} \sum_{i = 1}^{n}\frac{p^*(x_i - x_0;\theta_0, \epsilon)}{\hat{p}(x_i)}\left(y_i - (x_i - x_0)^{T}\beta - \alpha\right)^2
\end{align}
where $\hat{p}$ is some consistent estimate of $p$, the density of $x$. 

\begin{restatable}{asmpt}{} \label{assumption: lce and covariate density}
    We require that $p^*(x; \theta_0, \epsilon)$ be a zero-mean spherically symmetric distribution with isotropic covariance parameterized by $\theta_0$. Moreover, we require the following: 
    \begin{enumerate}
        \item (Compact support): the density $p^*(x; \theta_0, \epsilon)$ is bounded and satisfies $C(\theta_0, \epsilon) \subseteq \mathcal{B}(0, \epsilon)$ for \begin{equation} \label{eq: support}
             C(\theta_0, \epsilon) = \overline{\{x : p^*(x; \theta_0, \epsilon) \neq 0\}}.
        \end{equation}
        \item (LCE continuity): the density $p^*(x; \theta_0, \epsilon)$ is continuous in $x$ Lebesgue almost everywhere.
        \item (LCE covariance): the density $p^*(x; \theta_0, \epsilon)$ satisfies $\mathbb{E}_{p^*}[XX^T] \succ 0$. 
        \item (Covariate support): for all $x \in C(\theta_0, \epsilon)$ defined above, we have that $p_0(x) > m$ for some $m > 0$. Moreover, we assume that the support of the covariates is bounded, $\{x : p_0(x) \neq 0\} \subseteq \mathcal{B}(0, r). $
        \item (Covariate density smoothness): 
        there exists a non-negative integer $k$ such that the density $p_0(x)$ is $C^k$ and has bounded derivatives on $\mathbb{R}^d$.
    \end{enumerate} 
\end{restatable}
\noindent By a slight abuse of notation, we'll redefine each $x_i$ to have $x_0$ subtracted off and we'll let $X_i = \begin{bmatrix}
    1 & x_i
\end{bmatrix}^T$ to handle the intercept. The optimization becomes 
\begin{align} 
\widehat{\beta}_{x_0,\epsilon} &= \arg\min_{\beta \in \mathbb{R}^{d+1}} \sum_{i = 1}^{n}\underbrace{\frac{p^*(x_i;\theta_0, \epsilon)}{\hat{p}_0(x_i)}}_{\hat{w}_i}\left(y_i - X_i^{T}\beta\right)^2
\end{align}
where $\hat{p}_0$ is now a consistent estimate of the density of $x - \,x_0$. 
% Define the weighted means of $X_i$ and $Y_i$ as 
% \begin{equation}
%     \overline{X}_{\hat{w}} = \frac{\sum_{i = 1}^n \hat{w}(X_i)X_i}{\sum_{i = 1}^n \hat{w}(X_i)} \qquad \text{and} \qquad \overline{Y}_{\hat{w}} = \frac{\sum_{i = 1}^n \hat{w}(X_i)f(\Pi(X_i+x_0))}{\sum_{i = 1}^n \hat{w}(X_i)}.
% \end{equation}
Through a straightforward calculation one would find that
\begin{align}
    \widehat{\beta}_{x_0,\epsilon} &= \left(\sum_{i = 1}^n\hat{w}_iX_iX_i^T\right)^{-1}\left(\sum_{i = 1}^n\hat{w}_iX_iy_i\right). \label{eq: sample est closed form}
\end{align}
A similar set of calculations leads to the population quantity. We will again abuse notation and assume that $x$ has had $x_0$ subtracted off. The population optimization problem then looks like
\begin{align}
\beta^*_{x_0, \epsilon} &= \arg\min_{\beta \in \mathbb{R}^{d+1}} \mathbb{E}_{x\sim p_0} \Big[\frac{p^*(x;\theta_0, \epsilon)}{p_0(x)}\Big(f(\Pi(x + x_0)) - X^{T}\beta \Big)^2\Big] \label{eq: population estimator}
\end{align}
where $p_0$ is the density of the random variable $x$. Solving for $\beta^*_{x_0, \epsilon}$ yields
\begin{align}
    \beta^*_{x_0, \epsilon} = \mathbb{E}_{x \sim p^*(\cdot\,; \, \theta_0)}\left[XX^T\right]^{-1}\mathbb{E}_{x \sim p^*(\cdot\,;\,\theta_0)}\left[Xy\right] \label{eq: population est closed form}
\end{align}
where $y = f(\Pi(x+x_0))$. We will show that the last $d$ entries of $\widehat{\beta}_{x_0,\epsilon}$ are a consistent estimate of $\gamma'(\Pi(x_0))$ up to scale. To formalize this, let $V = \begin{bmatrix} 0 & I_d\end{bmatrix}$ and define the population approximation of $\gamma'(\Pi(x_0))$ as
\begin{equation}
    v^*_{x_0, \epsilon} = V\beta^*_{x_0, \epsilon} \label{eq: population velocity estimator}
\end{equation}
and the sample estimator
\begin{equation}
    \hat{v}_{x_0, \epsilon} = V\hat{\beta}_{x_0, \epsilon}. \label{eq: sample velocity estimator}
\end{equation}

\subsection{Density Estimation and Kernel Regression}
The estimators in \cref{eq: sample est closed form} and \cref{eq: sample velocity estimator} require an estimate of the density ratio $p^*(x_i;\theta_0, \epsilon)/p_0(x_i)$ -  we will need to impose some assumptions on a kernel density estimator $\hat{p}_0$ of $p_0$ that ensures consistency of the estimate of $\hat{v}_{x_0, \epsilon}$. Let $K(\cdot)$ be a kernel that meets \cref{assumption: kernel smoothness} and $h > 0$ its bandwidth parameter. Then the kernel density estimate of $p_0$ is
\begin{equation}
    \hat{p}_0(x) = \frac{1}{n}\sum_{i = 1}^nh^{-d}K\left(\frac{x -x_i}{h}\right). \label{eq: kde}
\end{equation}

\begin{restatable}[Kernel smoothness]{asmpt}{} \label{assumption: kernel smoothness}
    The kernel $K$ is Lipschitz, zero outside of a bounded set, and there exists a positive integer $s \leq k$ (with $k$ defined in \cref{assumption: lce and covariate density}) such that for all $d$-tuples of non-negative integers $(j_1, \dots, j_d)$ such that $\sum_i j_i < s$ and
    \[\int_{\mathbb{R}^d}u_1^{j_1}\dots u_d^{j_d}K(u)\,du_1\dots du_d.\]
\end{restatable}

\begin{restatable}[Boundedness and integrability, \cite{hansen2008uniform}]{asmpt}{} \label{assumption: boundedness and integrability}
    The kernel $K$ satisfies $\|K\|_{\infty} < \infty$ and 
    \[\int_{\mathbb{R}^d}|K(u)|\,du_1\dots du_d < \infty.\]
    Moreover, we require
    \[\int_{\mathbb{R}^d}K(u)\,du_1\dots du_d =  1.\]
\end{restatable}

\begin{restatable}[Uniform convergence over a compact set, \citet{87881b06-b1f4-38f8-874e-03ce94e584ba}]{lemma}{} \label{lemma: uniform convergence of kde}
    For a given covariate density $p_0(\cdot)$, let $C(\theta_0, \epsilon)$ be the compact set arising from \cref{assumption: lce and covariate density}. Suppose the kernel $K(\cdot)$ satisfies \cref{assumption: kernel smoothness} with integer $s$, and the bandwidth sequence $h_n$ is chosen such that $\sqrt{n}h_n^d/\ln(n) \rightarrow \infty$ and $\sqrt{n}h_n^s \rightarrow 0$. Then we have
    \begin{equation}
        \sup_{x \in C(\theta_0, \epsilon)} \left|\hat{p}_0(x) - p_0(x)\right| \overset{p}{\rightarrow} 0
    \end{equation}
    for $\hat{p}_0(\cdot)$ defined as in \cref{eq: kde}.
\end{restatable}

For the nonparametric regression kernel $K_r$ we will impose the same set of conditions, but add the following. 

\begin{restatable}[Local non-degeneracy of $K_r$]{asmpt}{} \label{assumption: kernel nondegeneracy}
    We assume that the nonparametric regression kernel $K_r$ satisfies \cref{assumption: kernel smoothness} and \cref{assumption: boundedness and integrability}. We also require that there exist constants $\delta > 0$ and $k_r > 0$ such that 
    \[|u| \leq \delta \implies K_r(u) \geq k_r.\]
\end{restatable}

\section{Consistency}

\subsection{Velocity estimate}

\begin{restatable}[Population quantity proportionality]{thm}{}\label{thm: population prop}
    For some scalar $\alpha_1 \in \mathbb{R}$ and Lebesgue almost every $x_0 \in \mathcal{B}(0, r)$ there exists an $\epsilon_0 > 0$ such that for all $\epsilon \leq \epsilon_0$ 
    \[\|v^*_{x_0, \epsilon} - \alpha_1\gamma'(s_0)\|_2 \lesssim (M_2+rM_3)L_f\epsilon\]
    with $v^*_{x_0, \epsilon}$ defined in \cref{eq: population velocity estimator} and $s_0 = \Pi(x_0)$.
\end{restatable}

\noindent \textit{Proof.} By \Cref{thm: differentiability of pi} we have that for Lebesgue almost every $x_0$, $\Pi(x_0)$ exists, is unique and $\nabla\Pi(x_0) \propto \gamma'(\Pi(x_0))$. Thus it suffices to show that $\beta^*_{x_0, \epsilon} \rightarrow \nabla\Pi(x_0)$. To do so, we'll define 
\[g(x) = f\Big(\Pi(x_0) + \nabla\Pi(x_0)^{T}(x-x_0)\Big)\]
and 
\[\xi(x) = f(\Pi(x)) - g(x).\] Note that \Cref{lemma: proj approx} implies that there exists an $\epsilon_0 > 0$ such that for all $\epsilon \leq \epsilon_0$ and for all $x \in \mathcal{B}(x_0, \epsilon)$ we have $\xi(x) \lesssim (M_2 + rM_3)L_f\epsilon^2$. Let  
\begin{align}
    \tilde{\beta}^*_{x_0, \epsilon} &= \arg\min_{\beta}\mathbb{E}_{x\sim p_0}\Big[\underbrace{\frac{p^*(x;\theta_0, \epsilon)}{p_{0}(x)}}_{w}\Big(g(x+x_0) - X^{T}\beta\Big)^2\Big]. \label{eq: perturbed pop estimator}
\end{align}
By \Cref{lemma: surrogate problem} we know that for some $\alpha_0, \alpha_1 \in \mathbb{R}$ we have $\tilde{\beta}^*_{x_0, \epsilon}  =  \begin{bmatrix}
    \alpha_0 & \alpha_1\nabla\Pi(x_0)
\end{bmatrix}^T$, and by \Cref{lemma: proj approx}, we know that
\begin{align*}
    \left|\xi(x+x_0)\right| &= \big|f(\Pi(\underbrace{x+x_0}_{z})) - f\left(\Pi(x_0) + \nabla\Pi(x_0)^Tx\right)\big| \\
    &= \left|f(\Pi(z)) - f\left(\Pi(x_0) + \nabla\Pi(x_0)^T(z-x_0)\right)\right|\\
    &\lesssim (M_2 + rM_3)L_f\epsilon^2.
\end{align*}
Observe that $\mathbb{E}_{x\sim p_{0}}[wXX^{T}]$ is necessarily positive definite as shown in \Cref{lemma: pd}. Thus, \Cref{lemma: perturbation} allows us to say 
\begin{align*} 
\|\beta^*_{x_0, \epsilon} - \tilde{\beta}_{x_0, \epsilon}^*\|_2 &= \left\|\left\{\mathbb{E}_{p_{0}}[wXX^{T}]\right\}^{-1}\mathbb{E}_{p_{0}}[wX\xi(x+x_0)]\right\|_2 \\
&\leq \left\|\left\{\mathbb{E}_{p_{0}}[wXX^{T}]\right\}^{-1}\mathbb{E}_{p_{0}}[wX]\right\|_2(M_2 + rM_3)L_f\epsilon^2.
\end{align*}
Now we'll apply \Cref{lemma: pnorm bound} to obtain,
\begin{align}
    \|\beta^*_{x_0, \epsilon} - \tilde{\beta}_{x_0, \epsilon}^*\|_2 &\lesssim (M_2 + rM_3)L_f\epsilon.
\end{align}
The statement follows from the definitions in \cref{eq: sample velocity estimator} and \cref{eq: population velocity estimator}. 

\qed{}

\begin{restatable}[Central limit for $\hat{v}_{x_0, \epsilon}$]{corollary}{} \label{corr: central limit for v}
    If \cref{assumption: kernel smoothness} and \cref{assumption: lce and covariate density} are satisfied, then we have asymptotic normality of the velocity vector estimate $\hat{v}_{x_0, \epsilon}$. In particular,
    \[\sqrt{n}(\hat{v}_{x_0, \epsilon} - v^*_{x_0, \epsilon}) \overset{d}{\rightarrow}N(0, VQ^{-1}\Lambda Q^{-1}V^T)\]
    with $V = \begin{bmatrix}
        0 & I_d
    \end{bmatrix}$, 
    \[\Lambda = \text{Var}(wXu)\]
    and and $u = y - X^T\beta_{x_0, \epsilon}^*$.
\end{restatable}
\noindent \textit{Proof.} The result follows from the central limit for $\hat{\beta}_{x_0, \epsilon}$ via \Cref{thm: central limit for beta} and application of the delta method with the Jacobian of the map $\hat{v}_{x_0, \epsilon} = V\hat{\beta}_{x_0, \epsilon}$. 

\qed{}

\begin{restatable}[Central limit for $\hat{\beta}_{x_0, \epsilon}$]{thm}{} \label{thm: decomposition for beta}
    If \cref{assumption: kernel smoothness} and \cref{assumption: lce and covariate density} are satisfied, then we have
    \[\sqrt{n}\left(\hat{\beta}_{x_0, \epsilon} - \beta^*_{x_0, \epsilon}\right) = \frac{1}{\sqrt{n}}Q^{-1}\sum_{i = 1}^nw_iX_iu_i +O_p\left(n^{\frac{-2s+d}{4s+2d}}\right)\]
    where 
    \[\Lambda = \text{Var}(wXu)\]
    and $u = y - X^T\beta_{x_0, \epsilon}^*$.
\end{restatable}
\noindent \textit{Proof.} This proof will follow closely the argument from \citet{Newey_Ruud_2005}. Let's define the matrices
\begin{equation}
    Q = \mathbb{E}_{p^*(\cdot\,;\, \theta_0)}\left[XX^T\right] \qquad \text{and} \qquad \hat{Q} = \frac{1}{n}\sum_{i = 1}^n\hat{w}_iX_iX_i^T.
\end{equation}
We have
\begin{align}
    \hat{\beta}_{x_0, \epsilon} - \beta^*_{x_0, \epsilon} &= \hat{Q}^{-1}\left(\frac{1}{n}\sum_{i = 1}^n\hat{w}_iX_iy_i\right) - \beta^*_{x_0, \epsilon} \\
    &= \hat{Q}^{-1}\left(\frac{1}{n}\sum_{i = 1}^n\hat{w}_iX_iy_i - \hat{Q}\beta^*_{x_0, \epsilon}\right) \\
    &= \hat{Q}^{-1}\left(\frac{1}{n}\sum_{i = 1}^n\hat{w}_iX_i\underbrace{\left(y_i - X_i^T\beta^*_{x_0, \epsilon}\right)}_{u_i}\right) \\
    &= \hat{Q}^{-1}\left(\frac{1}{n}\sum_{i = 1}^n\hat{w}_iX_iu_i\right).
\end{align}
Now we'll add and subtract out a term with the true weights $w_i$ to obtain
\begin{align}
    \hat{\beta}_{x_0, \epsilon} - \beta^*_{x_0, \epsilon} &= \hat{Q}^{-1}\left(\frac{1}{n}\sum_{i = 1}^nw_iX_iu_i + \frac{1}{n}\sum_{i = 1}^n(\hat{w}_i - w_i)X_iu_i\right) \\
    &= \hat{Q}^{-1}\left(\frac{1}{n}\sum_{i = 1}^n\frac{p^*(x_i; \theta_0, \epsilon)X_iu_i}{p_0(x_i)} + \frac{1}{n}\sum_{i = 1}^n\left[\frac{1}{\hat{p}_0(x_i)} - \frac{1}{p_0(x_i)}\right]p^*(x_i; \theta_0, \epsilon)X_iu_i\right).
\end{align}
Now if we define $a(z) = p^*(x_i; \theta_0, \epsilon)X_iu_i$ then we can write the above expression as
\begin{align}
    \hat{\beta}_{x_0, \epsilon} - \beta^*_{x_0, \epsilon} &= \hat{Q}^{-1}\left(\frac{1}{n}\sum_{i = 1}^n\frac{a(z)}{p_0(x_i)} + \frac{1}{n}\sum_{i = 1}^n\left[\frac{1}{\hat{p}_0(x_i)} - \frac{1}{p_0(x_i)}\right]a(z)\right).
\end{align}
Critically, $a(z)$ is supported on $\mathcal{B}(0, \epsilon)$ as indicated by \cref{assumption: lce and covariate density} -- we'll use this fact shortly. For now, we'll apply a first order Taylor expansion where $\Delta(x_i) = \hat{p}_0(x_i) - p_0(x_i)$ to say,
\begin{align}
    \hat{\beta}_{x_0, \epsilon} - \beta^*_{x_0, \epsilon} &= \hat{Q}^{-1}\bigg(\underbrace{\frac{1}{n}\sum_{i = 1}^n\frac{a(z)}{p_0(x_i)}}_{(I)} - \underbrace{\frac{1}{n}\sum_{i = 1}^n\frac{a(z)\Delta(x_i)}{p_0^2(x_i)}}_{(II)} + \underbrace{\frac{1}{n}\sum_{i = 1}^na(z) O(\Delta(x_i)^2)}_{(III)}\bigg).
\end{align}
Note that standard results by \citet{hansen2008uniform} indicate a uniform in $x$ (over $\mathcal{B}(0, \epsilon)$)  convergence of $\hat{p}_0$ to $p_0$ given by
\begin{equation}
    \|\Delta(x)\|_{\infty} = O_p\left(\sqrt{\frac{\log n}{nh^d}} + h^s\right)
\end{equation}
with $s$ from \cref{assumption: lce and covariate density}. With $h$ chosen optimally (in particular $h \asymp n^{-1/(2s+d)}$), we can achieve
\begin{equation}
    \|\Delta(x)\|_{\infty} = O_p\left(n^{\frac{-s}{2s+d}}\right)
\end{equation}
ignoring logarithmic factors. We will use this sup-norm bound soon, but for now note that $\mathbb{E}[a(z)/p_0(x)] = \mathbb{E}[wXu] = 0$ by the population first order conditions. We'll adopt the notation that $\mathbb{P}f = \mathbb{E}_{p_0}[f]$ and $\mathbb{P}_nf = \frac{1}{n}\sum_{i = 1}^nf(x_i)$. Thus,
\begin{align}
    (I) &= \frac{1}{n}\sum_{i = 1}^n\frac{a(z)}{p_0(x_i)} - \mathbb{E}\left[\frac{a(z)}{p_0(x)}\right] \\
    &= (\mathbb{P}_n -\mathbb{P})\left[\frac{a(z)}{p_0(x)}\right].
\end{align}
Since \Cref{lemma: Q convergence} establishes that $\hat{Q} \overset{p}\rightarrow Q$, we will introduce $Q$ and combine our new expression for $(I)$ as follows,
\begin{equation}
    \hat{\beta}_{x_0, \epsilon} - \beta^*_{x_0, \epsilon} = Q^{-1}\psi_n + (\hat{Q}^{-1} - Q^{-1})\psi_n.
\end{equation}
with
\begin{equation}
    \psi_n = (\mathbb{P}_n -\mathbb{P})\left[\frac{a(z)}{p_0(x)}\right] - \mathbb{P}_n\left[\frac{a(z)\Delta(x)}{p^2_0(x)}\right] + \mathbb{P}_n\left[a(z)O\left(\Delta(x)^2\right)\right].
\end{equation}
Now we'll consider the last three of the four terms of $\psi_n$ as they are of lower order. By \cref{assumption: lce and covariate density} we know that $\sup_{x\in \mathcal{B}(0, \epsilon)}\left|\frac{1}{p_0(x)}\right| \leq m^{-1}$ for some $m > 0$. Thus,
\begin{align*}
    \mathbb{P}_n\left[\frac{a(z)\Delta(x)}{p^2_0(x)}\right] &= \mathbb{P}_n\left[\frac{a(z)}{p_0(x)}\Delta(x)\right]\cdot O(1) \\
    &= \mathbb{P}_n\left[\frac{a(z)}{p_0(x)}\right]\cdot O_p\left(n^{\frac{-s}{2s+d}}\right) \\
    &= \mathbb{P}_n\left[wXu\right]\cdot O_p\left(n^{\frac{-s}{2s+d}}\right)\\
    &= o_p\left(n^{\frac{-s}{2s+d}}\right)
\end{align*}
where the last step follows from the LLN and the fact that $\mathbb{E}[wXu] = 0$ by population first order conditions. By the same logic, we have
\begin{align}
    \mathbb{P}_n\left[a(z)O\left(\Delta(x)^2\right)\right] &= O_p\left(n^{\frac{-2s}{2s+d}}\right)
\end{align}
Thus,
\begin{align}
    \psi_n = \frac{1}{n}\sum_{i = 1}^nw_iX_iu_i + O_p\left(n^{\frac{-2s}{2s+d}}\right).
\end{align}
Thus,
\begin{equation}
    \sqrt{n}\left(\hat{\beta}_{x_0, \epsilon} - \beta^*_{x_0, \epsilon}\right) = \frac{1}{\sqrt{n}}Q^{-1}\sum_{i = 1}^nw_iX_iu_i +O_p\left(n^{\frac{-2s+d}{4s+2d}}\right). 
\end{equation}

\qed{}

\begin{restatable}[Central limit for $\hat{\beta}_{x_0, \epsilon}$]{corollary}{} \label{thm: central limit for beta}
    If \cref{assumption: kernel smoothness} and \cref{assumption: lce and covariate density} are satisfied and $s > d/2$ then we have asymptotic normality of $\hat{\beta}_{x_0,\epsilon}$. In particular
    \[\sqrt{n}(\hat{\beta}_{x_0, \epsilon} - \beta^*_{x_0, \epsilon}) \overset{d}{\rightarrow} N(0, Q^{-1}\Lambda Q^{-1})\]
    where 
    \[\Lambda = \text{Var}(wXu)\]
    and $u = y - X^T\beta_{x_0, \epsilon}^*$.
\end{restatable}
\noindent\textit{Proof.} By the central limit theorem, the delta method, and \Cref{thm: decomposition for beta} we have
\begin{equation}
     \sqrt{n}(\hat{\beta}_{x_0, \epsilon} - \beta_{x_0, \epsilon}^*) \overset{p}{\rightarrow} N(0, Q^{-1}\Lambda Q^{-1})
\end{equation}
with $\Lambda = \text{Var}(wXu)$. 

\subsection{Nonparametric Regression}

To obtain a regression estimate for an $x$, we will project its neighborhood onto the estimate $\hat{v}_{x_0, \epsilon}$ and then perform kernel regression. Since we stipulate that $\gamma$ is parameterized by arclength, we will first normalize the velocity vector estimate to obtain $\hat{v}_{x_0, \epsilon}' = \hat{v}_{x_0, \epsilon}/\|\hat{v}_{x, \epsilon}\|_2$ (with its population counterpart being $v'^{*}_{x, \epsilon}$). 

\begin{algorithm}[H]
    \begin{algorithmic}[1]
        \caption{Locally Linear Single Index Regression}
        \label{alg: LLR}
        \Require New sample $x$, $\epsilon$-neighborhood of $x$ and associated responses $\{(x_i, y_i)\}_{i=1}^k \subseteq \mathbb{R}^d \times \mathbb{R}$, normalized velocity estimate $\hat{v}_{x, \epsilon}'$, a positive nonparametric regression kernel $K_r$ and bandwidth $h_r$\\
        Obtain neighborhood projection as $\hat{z}_i = \langle x_i - x, \hat{v}'_{x, \epsilon} \rangle$ \\
        Estimate response as 
        \begin{equation} \label{eq: kernel regressor}
            \hat{f}(x) = \sum_{i \in N(x, \epsilon)}\frac{y_iK_r(\frac{\hat{z}_i}{h_r})}{\sum_{j\in N(x, \epsilon)}K_r(\frac{\hat{z}_j}{h_r})}
        \end{equation}
        where $N(x, \epsilon) = \{i : \|x_i - x\|_2 \leq \epsilon\}$. \\
        \Return $\hat{f}(x)$
    \end{algorithmic}
\end{algorithm}

\begin{restatable}[Pointwise central limit of $\hat{f}$]{thm}{} \label{thm: regression ptwise clt}
    Let $x_1, \dots, x_n \sim \lambda$ where $x_i$ admits a density $p_0(x_i)$, and let $\hat{f}$ be an estimator of the regression function computed from the samples $x_1, \dots, x_n$ and their associated responses $y_1, \dots, y_n$. Under the assumptions of \Cref{corr: central limit for v} and for some $x\in \mathcal{B}(0, r)$ we have
    \[\sqrt{n}\left(\hat{f}(x) - f(s)\right) \overset{d}{\rightarrow} N(\bar{b}(x; \epsilon), \textcolor{red}{??}) \]
    where $s = \Pi(x)$ and $\bar{b}(x;\epsilon)$ is the asymptotic bias and has the form
    \begin{equation}
        \bar{b}(x; \epsilon) = \frac{\mathbb{E}_{x'}\left[\left\{f(\Pi(x')) - f(s)\right\}K_r(z'/h_r)\mathbbm{1}\{x'\in \mathcal{B}(x, \epsilon)\}\right]}{\mathbb{E}_{x'}\left[K_r(z'/h_r)\mathbbm{1}\{x'\in \mathcal{B}(x, \epsilon)\}\right]}
    \end{equation}
    with $z' = \langle x' - x, \gamma'(\Pi(x'))\rangle.$
\end{restatable}
\noindent \textit{Proof.} Let $z_i = \langle x_i - x, v'^*_{x, \epsilon}\rangle$ and define
\[\tilde{f}(x) = \sum_{i \in N(x, \epsilon)}\frac{y_iK_r(\frac{z_i}{h_r})}{\sum_{j \in N(x, \epsilon)}K_r(\frac{z_j}{h_r})}\]
and
\[\bar{f}(x) = \frac{\mathbb{E}_{x'}\left[f(\Pi(x'))K_r(z'/h_r)\mathbbm{1}\{x'\in \mathcal{B}(x, \epsilon)\}\right]}{\mathbb{E}_{x'}\left[K_r(z'/h_r)\mathbbm{1}\{x'\in \mathcal{B}(x, \epsilon)\}\right]}\]
where $z' = \langle x'-x,\gamma'(\Pi(x'))\rangle$. Expanding yields
\begin{align}
    \hat{f}(x) - f(x) &= \hat{f}(x) - \tilde{f}(x) + 
    \tilde{f}(x) - \bar{f}(x) + \bar{f}(x) - f(s).
\end{align}
Defining $A^{(1)}, B^{(1)}, \widehat{A}^{(1)}, \widehat{B}^{(1)}$ to be the numerator and denominator of $\tilde{f}$ and $\hat{f}$ respectively allows us to write the difference as
\begin{align}
    \hat{f}(x) - \tilde{f}(x)  &= \frac{\widehat{A}^{(1)}}{\widehat{B}^{(1)}} - \frac{A^{(1)}}{B^{(1)}} \\
    &= \frac{\widehat{A}^{(1)}-A^{(1)}}{B^{(1)}}\frac{B^{(1)}}{\widehat{B}^{(1)}} + \frac{A^{(1)}}{B^{(1)}}\left(\frac{B^{(1)} - \widehat{B}^{(1)}}{B^{(1)}}\right)\frac{B^{(1)}}{\widehat{B}^{(1)}} \\
    &= 
    \left[\frac{\widehat{A}^{(1)}-A^{(1)}}{B^{(1)}} + \frac{\frac{A^{(1)}}{B^{(1)}}\left(B^{(1)} - \widehat{B}^{(1)}\right)}{B^{(1)}}\right]\frac{B^{(1)}}{\widehat{B}^{(1)}}. \label{eq: kdr estimator decomp}
\end{align}
Now we can analyze these terms individually. Note that \Cref{lemma: regression ratios} indicates that if $\epsilon \lesssim h_r$ then
\begin{align*}
    \frac{B^{(1)}}{\widehat{B}^{(1)}} &= 1  + O_p\left(\frac{1}{h_r\sqrt{n}}\right) \\
    &= 1  + o_p\left(n^{-(d+2)/(2d+8)}\right)
\end{align*}
when $\epsilon \lesssim h_r$ and $\epsilon \asymp n^{-1/(d+4)}$. Now consider,
\begin{align}
    \frac{\widehat{A}^{(1)} - A^{(1)}}{B^{(1)}} &= \frac{\sum_{i=1}^n\mathbbm{1}\{x_i \in \mathcal{B}(x, \epsilon)\}y_i\left[K_r(\frac{\hat{z}_i}{h_r}) - K_r(\frac{z_i}{h_r})\right]}{\sum_{i = 1}^n\mathbbm{1}\{x_i \in \mathcal{B}(x, \epsilon)\}K_r(z_i/h_r)} \\
    &= \frac{\sum_{i=1}^n\mathbbm{1}\{x_i \in \mathcal{B}(x, \epsilon)\}y_i\left[\frac{d}{dz_i}K_r(\frac{z_i}{h_r})(\hat{z}_i - z_i) + O_p((\hat{z}_i - z_i)^2)\right]}{\sum_{i = 1}^n\mathbbm{1}\{x_i \in \mathcal{B}(x, \epsilon)\}K_r(z_i/h_r)} \\
    &= \frac{\sum_{i=1}^n\mathbbm{1}\{x_i \in \mathcal{B}(x, \epsilon)\}y_i\left[\frac{d}{dz_i}K_r(\frac{z_i}{h_r})(\hat{z}_i - z_i) + O_p((\hat{z}_i - z_i)^2)\right]}{\sum_{i = 1}^n\mathbbm{1}\{x_i \in \mathcal{B}(x, \epsilon)\}K_r(z_i/h_r)}.
\end{align}
Observe that $\hat{z}_i - z_i = \langle x_i - x, \hat{v}'_{x, \epsilon} - v'^*_{x, \epsilon}\rangle$. By \Cref{lemma: scaled high prob bound for velocity} we have
\begin{align*}
    O_p((\hat{z}_i - z_i)^2) &= O_p(\|x_i-x\|_2^2\|\|\hat{v}'_{x, \epsilon} - v'^*_{x, \epsilon}\|^2) \\
    &= O_p\left(\epsilon^2\left(\frac{1}{\epsilon^{d+2}n} + \epsilon^2\right) \right) \\
    &= O_p(\epsilon^{-d}n^{-1} + \epsilon^4) \\
    &= n^{-4/(d+4)}
\end{align*}
when $\epsilon \asymp n^{-1/(d+4)}.$ Thus,
\begin{align}
    \frac{\widehat{A}^{(1)} - A^{(1)}}{B^{(1)}} &\asymp \frac{\frac{1}{n}\sum_{i = 1}\mathbbm{1}\{x_i \in \mathcal{B}(x, \epsilon)\}y_i\frac{d}{dz_i}K_r(\frac{z_i}{h_r})O_p(n^{-2/(d+4)})}{\frac{1}{n}\sum_{i = 1}^n\mathbbm{1}\{x_i \in \mathcal{B}(x, \epsilon)\}K_r(z_i/h_r)} \\
    &= O_p(n^{-2/(d+4)}).
\end{align}
Now we'll consider $(B^{(1)} - \widehat{B}^{(1)})/B^{(1)}$ by applying the same Taylor expansion as above,
\begin{align}
    \frac{B^{(1)} - \widehat{B}^{(1)}}{B^{(1)}} &= \frac{\sum_{i = 1}^n\mathbbm{1}\{x_i \in \mathcal{B}(x, \epsilon)\}\left[K_r(z_i/h_r) - K_r(\hat{z}_i'/h_r)\right]
    }{\sum_{i = 1}^n\mathbbm{1}\{x_i \in \mathcal{B}(x, \epsilon)\}K_r(z_i/h_r)} \\
    &= -\frac{\sum_{i = 1}^n\mathbbm{1}\{x_i \in \mathcal{B}(x, \epsilon)\}\left[\frac{d}{dz_i}K_r\left(z_i/h_r\right)(\hat{z}_i - z_i) + O_p((\hat{z}_i - z_i)^2)\right]
    }{\sum_{i = 1}^n\mathbbm{1}\{x_i \in \mathcal{B}(x, \epsilon)\}K_r(z_i/h_r)} \\
    &\asymp -\frac{\sum_{i = 1}^n\mathbbm{1}\{x_i \in \mathcal{B}(x, \epsilon)\}\frac{d}{dz_i}K_r\left(z_i/h_r\right)(\hat{z}_i - z_i)
    }{\sum_{i = 1}^n\mathbbm{1}\{x_i \in \mathcal{B}(x, \epsilon)\}K_r(z_i/h_r)} \\
    &= -\frac{\frac{1}{n}\sum_{i = 1}^n\mathbbm{1}\{x_i \in \mathcal{B}(x, \epsilon)\}\frac{d}{dz_i}K_r\left(z_i/h_r\right)(\hat{z}_i - z_i)
    }{\frac{1}{n}\sum_{i = 1}^n\mathbbm{1}\{x_i \in \mathcal{B}(x, \epsilon)\}K_r(z_i/h_r)} \\
    &= O_p(n^{-2/(d+4)}).
\end{align}
Thus, plugging into \cref{eq: kdr estimator decomp} yields
\begin{equation}
    \hat{f}(x) - \tilde{f}(x) = O_p(n^{-2/(d+4)}).
\end{equation}
Now we'll apply a similar decomposition for $\tilde{f}(x) - \bar{f}(x)$. In particular, if we let $A^{(2)}$, $B^{(2)}$, $n\widehat{A}^{(2)}$, $n\widehat{B}^{(2)}$ be the numerator and denominator of $\bar{f}$ and $\tilde{f}$ respectively, then we have
\begin{align}
    \tilde{f}(x) - \bar{f}(x) &=  \left[\frac{\widehat{A}^{(2)}-A^{(2)}}{B^{(2)}} + \frac{\frac{A^{(2)}}{B^{(2)}}\left(B^{(2)} - \widehat{B}^{(2)}\right)}{B^{(2)}}\right]\frac{B^{(2)}}{\widehat{B}^{(2)}}. 
\end{align}
Now, by \Cref{lemma: scaled high prob bound for velocity} we have
\begin{align}
    \frac{B^{(2)}}{\widehat{B}^{(2)}} &= 1 + O_p\left(\frac{1}{\epsilon^{d/2}\sqrt{n}}\right) \\
    &= 1+ O_p\left(n^{-2/(d+4)}\right)
\end{align}
when $\epsilon \asymp n^{-1/(d+4)}.$ Now we can plug into the expression above to say
\begin{align}
    \tilde{f}(x) - \bar{f}(x) &= \left(1 + O_p(n^{-2/(d+4)})\right)\frac{1}{n}\sum_{i = 1}^n\phi(x_i, y_i; x, \epsilon) \\
    &= \frac{1}{n}\sum_{i = 1}^n\phi(x_i, y_i; x, \epsilon) + o_p(n^{-2/(d+4)})
\end{align}
where 
\begin{equation*}
    \phi(x_i, y_i; x, \epsilon, h_r) = \frac{M_i - \mathbb{E}M_i}{\mathbb{E}_{x'}\left[K_r(z'/h_r)\mathbbm{1}\{x'\in \mathcal{B}(x, \epsilon)\}\right]}
\end{equation*}
and 
\begin{equation}
    M_i = \mathbbm{1}\{x_i \in \mathcal{B}(x, \epsilon)\}K_r\left(\frac{z_i}{h_r}\right)\left[y_i - \bar{f}(x)\right].
\end{equation}
Observe that $\phi(x_i, y_i; x, \epsilon, h_r)$ is mean $0$. 
Trivially, for the third pair of terms we have
\begin{align}
    \bar{f}(x) - f(s) &= \frac{\mathbb{E}_{x'}\left[f(\Pi(x'))K_r(z'/h_r)\mathbbm{1}\{x'\in \mathcal{B}(x, \epsilon)\}\right]}{\mathbb{E}_{x'}\left[K_r(z'/h_r)\mathbbm{1}\{x'\in \mathcal{B}(x, \epsilon)\}\right]} - f(s) \\
    &= \frac{\mathbb{E}_{x'}\left[\left\{f(\Pi(x')) - f(s)\right\}K_r(z'/h_r)\mathbbm{1}\{x'\in \mathcal{B}(x, \epsilon)\}\right]}{\mathbb{E}_{x'}\left[K_r(z'/h_r)\mathbbm{1}\{x'\in \mathcal{B}(x, \epsilon)\}\right]}\\
    &= \bar{b}(x; \epsilon, h_r).
\end{align}
Putting it together, we have 
\begin{equation}
    \hat{f}(x) - f(s) = \bar{b}(x; \epsilon) + \frac{1}{n}\sum_{i = 1}^n\phi (x_i, y_i; x, \epsilon, h_r) + o_p(n^{-2/(d+4)})
\end{equation}
and thus,
\begin{equation}
    \sqrt{n}(\hat{f}(x) - f(s)) = \bar{b}(x; \epsilon, h_r) + \frac{1}{\sqrt{n}}\sum_{i = 1}^n\phi (x_i, y_i; x, \epsilon, h_r) + O_p(n^{-2/(d+4)})
\end{equation}
\textcolor{red}{need to verify the variance of $\phi$.}
\qed{}

\section{Risk}
\subsection{Velocity estimation}
\begin{restatable}[Risk of $\hat{v}_{x_0, \epsilon}$]{thm}{} \label{thm: velocity risk}
    Let $x_1, \dots, x_n \sim \lambda$ where $x_i$ admits a density $p_0(x_i)$. Under the assumptions of \Cref{corr: central limit for v} we have
    \[\|\hat{v}_{x_0, \epsilon} - \alpha_1\gamma'(s_0)\|_{L^2(\lambda)}^2 \lesssim \frac{1}{\epsilon^{d+2}n} + \epsilon^2 \]
    with $s_0 = \Pi(x_0)$. In particular, if we choose $\epsilon \asymp n^{-1/(d+4)}$ then we have
    \[\|\hat{v}_{x_0, \epsilon} - \alpha_1\gamma'(s_0)\|_{L^2(\lambda)}^2 \lesssim n^{-2/(d+4)}. \]
\end{restatable}
\noindent \textit{Proof.} By the triangle inequality and \Cref{thm: population prop} we have
\begin{align}
    \|\hat{v}_{x_0, \epsilon} - \alpha_1\gamma'(s_0)\|_{L^2(\lambda)}^2 &\leq \|\hat{v}_{x_0, \epsilon} - v^*_{x_0, \epsilon} \|_{L^2(\lambda)}^2 + \|v^*_{x_0, \epsilon}  - \alpha_1\gamma'(s_0)\|_{L^2(\lambda)}^2 \\
    &\lesssim  \|\hat{v}_{x_0, \epsilon} - v^*_{x_0, \epsilon} \|_{L^2(\lambda)}^2 + (M_2 + rM_3)^2L_f^2\epsilon^2.
\end{align}
To bound the first term, we'll use the fact that $\sqrt{n}(\hat{v}_{x_0, \epsilon} - v^*_{x_0, \epsilon}) \overset{d}{\rightarrow}N(0, VQ^{-1}\Lambda Q^{-1}V^T)$ from \Cref{corr: central limit for v} to say
\begin{align}
    \|\hat{v}_{x_0, \epsilon} - \alpha_1\gamma'(s_0)\|_{L^2(\lambda)}^2
    &\lesssim  \frac{1}{n}\text{tr}\left( VQ^{-1}\Lambda Q^{-1}V^T\right) + (M_2 + rM_3)^2L_f^2\epsilon^2.
\end{align}
By \Cref{lemma: trace bound} we have a bound on the trace of the asymptotic variance, leading to
\begin{align}
    \|\hat{v}_{x_0, \epsilon} - \alpha_1\gamma'(s_0)\|_{L^2(\lambda)}^2
    &\lesssim  \frac{1}{\epsilon^{d+2}n} + \epsilon^2.
\end{align}
If we choose $\epsilon \asymp n^{-1/(d+4)}$ then we have
\[\|\hat{v}_{x_0, \epsilon} - \alpha_1\gamma'(s_0)\|_{L^2(\lambda)}^2 \lesssim n^{-2/(d+4)}. \]

\qed{}

\subsection{Nonparametric Regression}

\begin{restatable}[Pointwise risk of $\hat{f}$]{thm}{} 
    Let $x_1, \dots, x_n \sim \lambda$ where $x_i$ admits a density $p_0(x_i)$. Under the assumptions of \Cref{corr: central limit for v} and for some $x\in \mathcal{B}(0, r)$ we have
    \[\|\hat{f}(x) - f(x)\|_{L^2(\lambda)}^2 \lesssim \textcolor{red}{something} \]
    with $s = \Pi(x)$, where $\hat{f}$ is an estimator of the regression function computed from the samples $x_1, \dots, x_n$.
\end{restatable}
\noindent \textit{Proof.} 

% Let $z_i = \Pi(x_i) - \Pi(x)$ and define
% \[\tilde{f}(x) = \sum_{i \in N(x, \epsilon)}\frac{y_iK_r(\frac{z_i}{h_r})}{\sum_{j \in N(x, \epsilon)}K_r(\frac{z_j}{h_r})}.\]
% Critically, observe that we are summing over indices $\{i : \|x_i - x\|_2 \leq \epsilon\}$, which is a subset of the full $n$ samples. Expanding yields
% \begin{align}
%     \|\hat{f}(x) - f(x) \|_{L^2(\lambda)}^2 &= \|\hat{f}(x) - \tilde{f}(x) + \tilde{f}(x) - f(s)\|_{L^2(\lambda)}^2 \\
%     &\leq \|\hat{f}(x) - \tilde{f}(x)\|_{L^2(\lambda)}^2 + \|\tilde{f}(x) - f(s)\|_{L^2(\lambda)}^2.
% \end{align}
% Note that the first term captures the projection error, while the second term captures the kernel regression error. Let's consider this first term. Defining $A, B, \widehat{A}, \widehat{B}$ to be the numerator and denominator of $\tilde{f}$ and $\hat{f}$ respectively allows us to write the difference as
% \begin{align}
%     \hat{f}(x) - \tilde{f}(x) &= \frac{\sum_{i \in N(x, \epsilon)}y_iK_r(\frac{\hat{z}_i}{h_r})}{\sum_{j \in N(x, \epsilon)}K_r(\frac{\hat{z}_j}{h_r})} - \frac{\sum_{i \in N(x, \epsilon)}y_iK_r(\frac{z_i}{h_r})}{\sum_{j \in N(x, \epsilon)}K_r(\frac{z_j}{h_r})} \\
%     &= \frac{\widehat{A}}{\widehat{B}} - \frac{A}{B} \\
%     &= \frac{\widehat{A}-A}{B}\frac{B}{\widehat{B}} + \frac{A}{B}\left(\frac{B - \widehat{B}}{B}\right)\frac{B}{\widehat{B}}.
% \end{align}
% Now we can analyze these terms individually. Note that 
% \begin{align*}
%     \frac{B}{\hat{B}} &= \frac{\sum_{j \in N(x, \epsilon)}K_r(\frac{z_j}{h_r})}{\sum_{j \in N(x, \epsilon)}K_r(\frac{\hat{z}_j}{h_r})} \\
%     &= \frac{\sum_{j \in N(x, \epsilon)}K_r(\frac{z_j}{h_r})}{\sum_{j \in N(x, \epsilon)}K_r(\frac{z_j + \hat{z}_j - z_j}{h_r})}
% \end{align*}

% Now define
% \[\bar{f}(x) = \sum_{i \in N(x, \epsilon)}\frac{y_iK_r(\frac{\hat{z}_i}{h_r})}{\sum_{j \in N(x, \epsilon)}K_r(\frac{z_j}{h_r})}\]
% which allows us to say
% \begin{align}
%     \|\hat{f}(x) - \tilde{f}(x)\|_{L^2(\lambda)}^2 &= \|\hat{f}(x) - \bar{f}(x) + \bar{f}(x) - \tilde{f}(x)\|_{L^2(\lambda)}^2 \\
%     &\leq \|\hat{f}(x) - \bar{f}(x)\|_{L^2(\lambda)}^2 + \| \bar{f}(x) - \tilde{f}(x)\|_{L^2(\lambda)}^2.
% \end{align}
% Expanding the first term,
% \begin{align}
%     \|\hat{f}(x) - \bar{f}(x)\|_{L^2(\lambda)}^2&= \left\|\frac{\sum_{i \in N(x, \epsilon)}y_i\left[K_r(\frac{\hat{z}_i}{h_r}) - K_r(\frac{z_i}{h_r})\right]}{\sum_{j \in N(x, \epsilon)}K_r(\frac{\hat{z}_j}{h_r})}\right\|_{L^2(\lambda)}^2.
% \end{align}
% We can bound the discrepancy in the numerator using the Lipschitz property of the kernel (as guaranteed by \Cref{assumption: kernel smoothness}) and \Cref{lemma: projection error}. In particular, we have
% \begin{align}
%     K_r(\hat{z}_i/h_r) - K_r(z_i/h_r) &\leq L_{K_r}|\hat{z}_i - z_i| \\
%     &\leq \epsilon L_{K_r}\left(\|\hat{v}_{x, \epsilon} - \gamma'(s)\|_2 + L_{x, \epsilon} + 1\right)
% \end{align}
% where $L_{K_r}$ is taken to be the Lipschitz constant of $K_r$. Thus,
% \begin{align}
%     \|\hat{f}(x) - \bar{f}(x)\|_{L^2(\lambda)}^2&\leq \left\|\frac{\sum_{i \in N(x, \epsilon)}y_i\left|K_r(\frac{\hat{z}_i}{h_r}) - K_r(\frac{z_i}{h_r})\right|}{\sum_{j \in N(x, \epsilon)}K_r(\frac{\hat{z}_j}{h_r})}\right\|_{L^2(\lambda)}^2 \\
%     &\leq \left\|\frac{\epsilon L_{K_r}\left(\|\hat{v}_{x, \epsilon} - \gamma'(s)\|_2 + L_{x, \epsilon} + 1\right)\sum_{i \in N(x, \epsilon)}y_i}{\sum_{j \in N(x, \epsilon)}K_r(\frac{\hat{z}_j}{h_r})}\right\|_{L^2(\lambda)}^2 \\
%     &= \epsilon L_{K_r}\left(\|\hat{v}_{x, \epsilon} - \gamma'(s)\|_2 + L_{x, \epsilon} + 1\right)\left\|\frac{\sum_{i \in N(x, \epsilon)}y_i}{\sum_{j \in N(x, \epsilon)}K_r(\frac{\hat{z}_j}{h_r})}\right\|_{L^2(\lambda)}^2
% \end{align}
% Expanding the second term,
% \begin{align*}
%     \| \bar{f}(x) - \tilde{f}(x)\|_{L^2(\lambda)}^2 &= \left\|\sum_{i \in N(x, \epsilon)}y_i K_r(z_i/h_r)\left[\frac{1}{\sum_{j \in N(x, \epsilon)}K_r(\frac{z_j}{h_r})} - \frac{1}{\sum_{j \in N(x, \epsilon)}K_r(\frac{\hat{z}_j}{h_r})}\right]\right\|_{L^2(\lambda)}^2
% \end{align*}

\section{Technical results}

\begin{restatable}{lemma}{} \label{lemma: scaled high prob bound for velocity}
    If \cref{assumption: kernel smoothness} and \cref{assumption: lce and covariate density} are satisfied, then we have 
    \[\|\hat{v}_{x_0, \epsilon} - v^*_{x_0, \epsilon}\|_2 = O_p\left(\frac{1}{\epsilon\sqrt{n}}\right).\]
    Moreover, if $\alpha_1 \neq 0$ then 
    \[\|\hat{v}'_{x_0, \epsilon} - v'^*_{x_0, \epsilon}\|_2 = O_p\left(\frac{1}{\epsilon\sqrt{n}}\right).\]
\end{restatable}
\noindent \textit{Proof.} By the delta method and \Cref{corr: central limit for v} we have
\[\epsilon\sqrt{n}(\hat{v}_{x_0, \epsilon} - v^*_{x_0, \epsilon}) \overset{d}\rightarrow N(0, \epsilon^2VQ^{-1}\Lambda Q^{-1}V^T). \]
Now by \Cref{lemma: trace bound} we know that $\text{tr}(\epsilon^2VQ^{-1}\Lambda Q^{-1}V^T) = O(1)$, thus
\[\epsilon\sqrt{n}\|\hat{v}_{x_0, \epsilon} - v^*_{x_0, \epsilon}\|_2 = O_p(1).\]
The first result follows. For the second result, note that by \Cref{thm: population prop} $\|v^*_{x_0, \epsilon}\|_2 \rightarrow \alpha_1 > 0$ as $\epsilon \rightarrow 0$. Thus, by the first result $\|\hat{v}_{x_0, \epsilon} - v^*_{x_0, \epsilon}\|_2 \leq \|v^*_{x_0, \epsilon}\|_2/2$ with probability $\rightarrow 1$. It follows that 
\begin{align}
    \|\hat{v}'_{x_0, \epsilon} - v'^*_{x_0, \epsilon}\|_2 &= \left\|\frac{\hat{v}_{x_0, \epsilon}}{\|\hat{v}_{x_0, \epsilon}\|_2} - \frac{v^*_{x_0, \epsilon}}{\|v^*_{x_0, \epsilon}\|_2}\right\|_2 \\
    &= \frac{1}{\|v^*_{x_0, \epsilon}\|_2}\left\|\hat{v}_{x_0, \epsilon}\frac{\|v^*_{x_0, \epsilon}\|_2}{\|\hat{v}_{x_0, \epsilon}\|_2} - v^*_{x_0, \epsilon}\right\|_2.
\end{align}
Now we'll use the first result to bound the ratio inside the norm. In particular, by triangle inequality
\begin{align}
    \left|\|\hat{v}_{x_0, \epsilon}\|_2 - \|v^*_{x_0, \epsilon}\|_2\right| &\leq \|\hat{v}_{x_0, \epsilon} - v^*_{x_0, \epsilon}\|_2 \\
    &= O_p\left(\frac{1}{\epsilon\sqrt{n}}\right).
\end{align}
Thus,
\begin{align}
    \|\hat{v}_{x_0, \epsilon}\|_2 &= \|v^*_{x_0, \epsilon}\|_2 + O_p\left(\frac{1}{\epsilon\sqrt{n}}\right)
\end{align}
from which it follows
\begin{align}
    \frac{\|v^*_{x_0, \epsilon}\|_2}{\|\hat{v}_{x_0, \epsilon}\|_2} &= 1 + O_p\left(\frac{1}{\epsilon\sqrt{n}}\right)
\end{align}
since $\alpha_1 \neq 0$. Plugging back in yields
\begin{align}
    \|\hat{v}'_{x_0, \epsilon} - v'^*_{x_0, \epsilon}\|_2 &=  \frac{1}{\|v^*_{x_0, \epsilon}\|_2}\left\|\hat{v}_{x_0, \epsilon} - v^*_{x_0, \epsilon} + O_p\left(\frac{\|\hat{v}_{x_0, \epsilon}\|}{\epsilon\sqrt{n}}\right)\right\|_2 \\
    &\leq \frac{1}{\|v^*_{x_0, \epsilon}\|_2}\left\|\hat{v}_{x_0, \epsilon} - v^*_{x_0, \epsilon} \right\|_2 + O_p\left(\frac{\|\hat{v}_{x_0, \epsilon}\|}{\|v^*_{x_0, \epsilon}\|_2\epsilon\sqrt{n}}\right) \\
    &= O_p\left(\frac{1}{\epsilon\sqrt{n}}\right) .
\end{align}

\qed{}

\begin{restatable}{lemma}{} \label{lemma: regression ratios}
    Let $\widehat{B}^{(1)}$, $B^{(1)}$, $\widehat{B}^{(2)}$ and $B^{(2)}$ be defined as in \Cref{thm: regression ptwise clt}. If we have $\epsilon \lesssim  h_r$, then
    \begin{equation}
        \frac{B^{(1)}}{\widehat{B}^{(1)}} = 1 +  O_p\left(\frac{1}{h_r\sqrt{n}}\right)
    \end{equation}
    and 
    \begin{equation}
        \frac{B^{(2)}}{\widehat{B}^{(2)}} = 1 + O_p\left(\frac{1}{\epsilon^{d/2}\sqrt{n}}\right).
    \end{equation}
\end{restatable}
\noindent \textit{Proof.} We'll start with proving the first equation. We have
\begin{align}
    \left|\frac{B^{(1)}}{\widehat{B}^{(1)}} - 1\right| &= \frac{\left|B^{(1)} - \widehat{B}^{(1)}\right|}{\left|\widehat{B}^{(1)}\right|}.
\end{align}
Let $\delta_i = \hat{z}_i - z_i = \langle x_i - x, \hat{v}'_{x_0, \epsilon} - v'^*_{x_0, \epsilon} \rangle$. Observe that
\begin{align*}
    \left|\delta_i/h_r\right| &\leq \|x_i - x\|_2\|\hat{v}'_{x_0, \epsilon} - v'^*_{x_0, \epsilon}\|_2/h_r
\end{align*}
and therefore by Lipschitzness
\begin{align}
    \left|B^{(1)} - \widehat{B}^{(1)}\right| &= \left|\frac{1}{n}\sum_{j = 1}^{n}\mathbbm{1}\{x_j \in \mathcal{B}(x, \epsilon)\}\left(K_r(z_j/h_r) - K_r(\hat{z}_j/h_r)\right)\right| \\
    &\leq \left|\frac{\|\hat{v}'_{x_0, \epsilon} - v'^*_{x_0, \epsilon}\|_2L_{K_r}\epsilon}{h_r}\frac{1}{n}\sum_{j = 1}^{n}\mathbbm{1}\{x_j \in \mathcal{B}(x, \epsilon)\}\right| \\
    &= \frac{\|\hat{v}'_{x_0, \epsilon} - v'^*_{x_0, \epsilon}\|_2L_{K_r}\epsilon N_{\epsilon}}{nh_r}
\end{align}
where $N_\epsilon = \sum_{j = 1}^{n}\mathbbm{1}\{x_j \in \mathcal{B}(x, \epsilon)\}$. It follows that $N_{\epsilon}/n \asymp \mathbb{P}(x' \in \mathcal{B}(x,\epsilon)) \asymp \epsilon^d$. Also note that \Cref{lemma: scaled high prob bound for velocity} indicates that 
\[\|\hat{v}'_{x_0, \epsilon} - v'^*_{x_0, \epsilon}\|_2 = O_p\left(\frac{1}{\epsilon\sqrt{n}}\right).\]
Putting it together,
\[\left|B^{(1)} - \widehat{B}^{(1)}\right| = O_p\left(\frac{L_{K_r} \epsilon^d}{\sqrt{n}h_r}\right).\]
Now we'll consider the denominator. Observe that $\hat{z}_j \leq \epsilon$ for all $x_j \in \mathcal{B}(x, \epsilon)$. Now we will invoke \Cref{assumption: kernel nondegeneracy}, which guarantees a neighborhood of size $\delta$ about $0$ where $K_r$ is strictly positive and $K_r \geq k_r > 0$. In particular, if we ensure that $|\epsilon/h_r| \leq \delta$ then we guarantee that $K_r \geq k_r$. Thus, if $\epsilon/\delta h_r = o(1)$ then 
\[\widehat{B}^{(1)} = \frac{1}{n}\sum_{i = 1}^n\mathbbm{1}\{x_j \in \mathcal{B}(x, \epsilon)\}K_r(\hat{z}_i/h_r) \gtrsim \epsilon^d.\]
It follows that 
\begin{align}
    \left|\frac{B^{(1)}}{\widehat{B}^{(1)}} - 1\right| &= O_p\left(\frac{1}{\sqrt{n}h_r}\right).
\end{align}
This completes the first result. For the second result, we'll apply similar logic. In particular,
\begin{align}
    \left|\frac{B^{(2)}}{\widehat{B}^{(2)}} - 1\right| &= \frac{\left|B^{(2)} - \widehat{B}^{(2)}\right|}{\left|\widehat{B}^{(2)}\right|}.
\end{align}
For the numerator, we have 
\begin{align}
    \left|B^{(2)} - \widehat{B}^{(2)}\right| &= \left|\frac{1}{n}\sum_{i = 1}^n\mathbbm{1}\{x_i \in \mathcal{B}(x, \epsilon)\}K_r(z_i/h_r) - \mathbb{E}\left[\mathbbm{1}\{x' \in \mathcal{B}(x, \epsilon)\}K_r(z'/h_r)\right]\right| \\
    &= O_p\left(\sqrt{\text{Var}\left(\widehat{B}^{(2)} - B^{(2)}\right)}\right) \\
    &= O_p\left(\sqrt{\text{Var}\left(\widehat{B}^{(2)}\right)}\right) \\
    &= O_p\left(\sqrt{\frac{\text{Var}\left(\mathbbm{1}\{x_i \in \mathcal{B}(x, \epsilon)\}K_r(z_i/h_r)\right)}{n}}\right) \\
    &= O_p\left(\sqrt{\frac{\|K_r\|_{\infty}\mathbb{E}\left[\mathbbm{1}\{x_i \in \mathcal{B}(x, \epsilon)\}^2\right]}{n}}\right) \\
    &= O_p\left(\frac{\epsilon^{d/2}}{\sqrt{n}}\right).
\end{align}
Since we also have 
\[\widehat{B}^{(2)} = \frac{1}{n}\sum_{i = 1}^n\mathbbm{1}\{x_j \in \mathcal{B}(x, \epsilon)\}K_r(z_i/h_r) \gtrsim \epsilon^d\]
it follows that 
\begin{equation}
    \left|\frac{B^{(2)}}{\widehat{B}^{(2)}} - 1\right| = O_p\left(\frac{1}{\epsilon^{d/2}\sqrt{n}}\right).
\end{equation}
\qed{}

% \begin{restatable}[Projection error]{lemma}{} \label{lemma: projection error limit}
%     For a point $x\in R$ (with $R$ defined in \Cref{thm: differentiability of pi}) we have for any $x_i \in \mathcal{B}(x, \epsilon)$, $\langle x_i - x, \hat{v}_{x, \epsilon}'\rangle  = \langle x_i - x, v'^*_{x, \epsilon}\rangle + O_p( n^{-1/2})$. 
% \end{restatable}
% \noindent \textit{Proof.} From \Cref{corr: central limit for v} we know that  
% \[\alpha_1\sqrt{n}(\hat{v}'_{x, \epsilon} - v'^*_{x, \epsilon}) \overset{d}\rightarrow N(0, \Sigma(\epsilon)).\]
% Let $e_i = x_i - x$ and $\Delta v = \hat{v}_{x, \epsilon}' - v'^*_{x, \epsilon}$. Then by the delta method
% \begin{align}
%     \langle e_i, \Delta v \rangle \overset{d}\rightarrow N(0, n^{-1}\alpha_1^{-2}e_i^T\Sigma(\epsilon)e_i).
% \end{align}
% Thus, 
% \begin{align}
% \text{Var}(\langle e_i, \Delta v\rangle) &\asymp  \frac{1}{n\alpha_1^{2}}e_i^T\Sigma(\epsilon)e_i \\
% &\leq \frac{\epsilon^2}{n\alpha_1^{2}}\|\Sigma(\epsilon)\|_{\text{op}} \\
% &\leq \frac{\epsilon^2}{n\alpha_1^{2}}\text{tr}(\Sigma(\epsilon))
% \end{align}
% where the last step followed from the fact that $\Sigma(\epsilon) \succeq 0$. Due to \Cref{lemma: trace bound} we have a bound on the trace, yielding
% \begin{align}
%     \text{Var}(\langle e_i, \Delta v\rangle)  = O_p\left(\frac{1}{n\alpha_1^2}\right).
% \end{align}
% Thus,
% \begin{align}
%     \langle x_i - x, \hat{v}_{x, \epsilon}'\rangle  -  \langle x_i - x, v'^*_{x, \epsilon}\rangle &= O_p( n^{-1/2})
% \end{align}
% from which the statement follows.

% \qed{}


% \begin{restatable}[Projection error]{lemma}{} \label{lemma: projection error}
%     For a point $x\in R$ (with $R$ defined in \Cref{thm: differentiability of pi}) and fixed point cloud $x_1, \dots, x_n$ contained in the support of the measure $\lambda$ such that $\|x_i - x\|_2 \leq \epsilon$, we have
%     \begin{equation}
%         \left||\langle x_i - x, \hat{v}_{x_0, \epsilon}'\rangle| - |\Pi(x_i) - \Pi(x)| \right| \leq \left(\|\hat{v}_{x, \epsilon} - \gamma'(s)\|_2 + L_{x, \epsilon} + 1\right)\epsilon
%     \end{equation}
%     where $\hat{v}'_{x, \epsilon} = \hat{v}'_{x, \epsilon}(x_1, \dots, x_n, y_1, \dots, y_n)$ can be regarded as a function of $x_1, \dots, x_n$ and their associated responses $y_1, \dots, y_n$. 
% \end{restatable}
% \noindent \textit{Proof.} Letting $s = \Pi(x)$ and expanding the left-hand side yields
% \begin{align*}
%      \left||\langle x_i - x, \hat{v}_{x, \epsilon}'\rangle| - |\Pi(x_i) - \Pi(x)| \right| &=  \left|\,\left|\langle x_i - x, \hat{v}_{x, \epsilon}' - \gamma'(s) + \gamma'(s)\rangle\right| - \left|\Pi(x_i) - \Pi(x)\right|\, \right| \\
%      &= \left|\, \left|\langle x_i - x, \hat{v}_{x, \epsilon}' - \gamma'(s)\rangle + \langle x_i - x, \gamma'(s)\rangle\right| - \left|\Pi(x_i) - \Pi(x)\right| \,\right| \\
%      &\leq \left|\,\left|\epsilon\|\hat{v}_{x, \epsilon} - \gamma'(s)\|_2 + \epsilon\right| - |\Pi(x_i) - \Pi(x)| \,\right|.
% \end{align*}
% where the last step follows from application of Cauchy Schwartz. By the assumption that $x \in R$ we know that $\Pi$ is differentiable on $\mathcal{B}(0, \epsilon)$ for $\epsilon$ sufficiently small. Differentiability of $\Pi$ implies \textit{locally} Lipschitz. Taking $\mathcal{B}(0, \epsilon)$ to be the open ball, we know that $\Pi$ is Lipschitz on $\mathcal{B}(0, \epsilon)$ with some constant $L_{x, \epsilon}$ non-increasing in $\epsilon$. This allows us to bound $|\Pi(x_i) - \Pi(x)| \leq L_{x, \epsilon}\epsilon$, yielding
% \begin{align*}
%      \left||\langle x_i - x, \hat{v}_{x, \epsilon}'\rangle| - |\Pi(x_i) - \Pi(x)| \right| &\leq \left|\,\left|\epsilon\|\hat{v}_{x, \epsilon} - \gamma'(s)\|_2 + \epsilon\right| - L_{x,\epsilon}\epsilon \,\right| \\
%      &\leq \left(\|\hat{v}_{x, \epsilon} - \gamma'(s)\|_2 + L_{x, \epsilon} + 1\right)\epsilon.
% \end{align*}

% \qed{}

% \begin{restatable}[Projection error limit]{corollary}{} \label{corollary: projection error limit}
%     For a point $x\in R$ (with $R$ defined in \Cref{thm: differentiability of pi}) and taking $\epsilon \asymp n^{-1/4}$ yields $\left||\langle x_i - x, \hat{v}_{x_0, \epsilon}'\rangle| - |\Pi(x_i) - \Pi(x)| \right| \overset{p}{\rightarrow} 0$. 
% \end{restatable}
% \noindent \textit{Proof.} Combining \Cref{lemma: projection error} and \Cref{thm: velocity risk} yields
% \[\left||\langle x_i - x, \hat{v}_{x_0, \epsilon}'\rangle| - |\Pi(x_i) - \Pi(x)|\right| \lesssim (n^{-1/4} + O(1) + 1)n^{-1/4}\]
% when $\epsilon \asymp n^{-1/4}$. Note that $L_{x, \epsilon}$ is non-increasing in $\epsilon$, so we can replace it with $O(1)$ to yield the inequality above. The statement follows by taking $n \rightarrow \infty$. 

% \qed{}




\begin{restatable}{lemma}{} \label{lemma: trace bound}
    For the matrices defined in \Cref{thm: central limit for beta} and appropriately chosen $p^*$, we have
    \[\text{tr}(VQ^{-1}\Lambda Q^{-1}V^T) \lesssim \epsilon^{-d-2}. \]
\end{restatable}
\noindent \textit{Proof.} Observe that we can rewrite $Q$ as follows
\begin{align}
    Q &= \mathbb{E}_{p^*}\left[\begin{pmatrix}
        1 \\ x
    \end{pmatrix}\begin{pmatrix}
        1 \\ x
    \end{pmatrix}^T\right] \\
    &= \mathbb{E}_{p^*}\left[
    \begin{bmatrix}
         1 & x^T \\
         x & xx^T
    \end{bmatrix}\right] \\
    &= \begin{bmatrix}
         \mathbb{P}_{p^*}(x \in \mathcal{B}(0, \epsilon)) & 0 \\
         0 & \mathbb{E}\left[xx^T\right]
    \end{bmatrix} \\
    &= \begin{bmatrix}
         1 & 0 \\
         0 & \nu^2I_d
    \end{bmatrix}
\end{align}
where
\[\nu^2 = \mathbb{E}_{p^*}[x_i^2]\]
for any $i \in [d]$. Note that the last two steps follows from the spherical symmetry of $p^*$. It follows that $Q^{-1}$ is
\begin{align}
    Q^{-1} &= \begin{bmatrix}
         1 & 0 \\
         0 & \nu^{-2}I_d
    \end{bmatrix}.
\end{align}
Now we'll analyze $\Lambda$, 
\begin{align}
    \Lambda &= \mathbb{E}_{p_0}\left[w^2\text{Var}(u)\begin{pmatrix}
        1 \\ x
    \end{pmatrix}\begin{pmatrix}
        1 \\ x
    \end{pmatrix}^T\right] \\
    &= 
    \text{Var}(u)\cdot\begin{bmatrix}
         \mathbb{E}[w^2] & \mathbb{E}[w^2x] \\
         \mathbb{E}[w^2x] & \mathbb{E}[w^2xx^T]
    \end{bmatrix}.
\end{align}
From the definition of $V$ we know that applying it in the manner above will select the bottom right $d\times d$ entries of the block diagonal matrix $Q^{-1}\Lambda Q^{-1}$. Therefore, after performing the matrix multiplications one obtains
\begin{align}
    VQ^{-1}\Lambda Q^{-1}V^T&= \text{Var}(u)\nu^{-4} \cdot \mathbb{E}_{p_0}\left[w^2xx^T\right] \\
    &= \text{Var}(u)\nu^{-4} \cdot \mathbb{E}_{p^*}\left[wxx^T\right]
\end{align}
By linearity of trace,
\begin{align}
    \text{tr}(VQ^{-1}\Lambda Q^{-1}V^T) &= \text{Var}(u)\nu^{-4} \cdot \mathbb{E}_{p^*}\left[\text{tr}(wxx^T)\right] \\
    &= \text{Var}(u)\nu^{-4} \cdot \mathbb{E}_{p^*}\left[\sum_{i = 1}^dwx_i^2\right] \\
    &= d\text{Var}(u)\nu^{-4} \cdot \mathbb{E}_{p^*}\left[wx_i^2\right] \\
    &\leq w_{\max}d\text{Var}(u)\nu^{-4} \cdot \underbrace{\mathbb{E}_{p^*}\left[x_i^2\right]}_{\nu^2} \\
    &= w_{\max}d\text{Var}(u)\nu^{-2}
\end{align}
where $w_{\max} = \max_{x\in \mathcal{B}(0, \epsilon)}w.$ Since $\text{Var}(u)$ is the variance of bounded random variables, it is $O(1)$ with respect to $\epsilon$. Thus, if we choose $p^*$ such that $\nu^2 \asymp \epsilon^2$ and $w_{\max} \asymp \epsilon^{-d}$ then the statement follows. 

\qed{}

\begin{restatable}{lemma}{}
    Let the matrices $Q$ and $\hat{Q}$ be defined as follows, 
    \begin{equation}
    Q = \mathbb{E}_{p^*(\cdot\,;\, \theta_0)}\left[XX^T\right] \qquad \text{and} \qquad \hat{Q} = \frac{1}{n}\sum_{i = 1}^n\hat{w}_iX_iX_i^T.
    \end{equation}
    Under \cref{assumption: kernel smoothness} and \cref{assumption: lce and covariate density} we have 
    \[\left\|\hat{Q}^{-1} - Q^{-1}\right\|_{\text{op}} = O_p\left(n^{\frac{-s}{2s+d} \wedge -\frac{1}{2}}\right)\]
    and thus $\hat{Q}^{-1}\overset{p}{\rightarrow}Q^{-1}$. \label{lemma: Q convergence}
\end{restatable}
\noindent \textit{Proof.} Recall that $Q$ can be written as an expectation over $p_0$, $Q = \mathbb{E}_{p_0}\left[wXX^T\right]$. Before establishing the convergence of $\hat{Q}$, we will bound $\hat{w}_i - w_i$ as follows
\begin{align*}
    \hat{w}_i - w_i &= p^*(x_i; \theta_0, \epsilon) \left[\frac{1}{\hat{p}_0(x_i)} - \frac{1}{p_0(x_i)}\right] \\
    &= p^*(x_i; \theta_0, \epsilon) \left[\frac{\Delta(x_i)}{p^2_0(x_i)} + O(\Delta(x_i)^2)\right]
\end{align*}
where $\Delta(x_i) = p_0(x_i) - \hat{p}_0(x_i)$. From \citet{hansen2008uniform} we know that 
\begin{equation}
    \sup_{x \in \mathcal{B}(0, \epsilon)} |\Delta(x)| = O_p\left(\sqrt{\frac{\log n}{nh^d}} + h^s\right)
\end{equation}
where $s$ comes from \Cref{assumption: lce and covariate density} and $h$ is the bandwidth for the kernel density estimate $\hat{p}_0$. For $h$ chosen optimally, we have
\begin{equation}
    \max_{i \in \{1, \dots, n \}} |\hat{w}_i - w_i| = O_p\left(n^{\frac{-2s}{2s+d}}\right).
\end{equation}
Now define 
\[\tilde{Q} = \frac{1}{n}\sum_{i = 1}^n w_iX_iX_i^T\]
and observe that 
\begin{align}
\left\|\hat{Q} - \tilde{Q}\right\|_{\text{op}} &= \left\|\frac{1}{n}\sum_{i = 1}^n\hat{w}_iX_iX_i^T - \frac{1}{n}\sum_{i = 1}^nw_iX_iX_i^T\right\|_{\text{op}} \\
&= \left\|\frac{1}{n}\sum_{i = 1}^n(\hat{w}_i - w_i)X_iX_i^T\right\|_{\text{op}} \\
&\leq \max_{i \in \{1, \dots, n\}}|\hat{w}_i - w_i| \left\|\frac{1}{n}\sum_{i = 1}^nX_iX_i^T\right\|_{\text{op}} \\
&= O_p\left(n^{\frac{-2s}{2s+d}}\right).
\end{align}
By triangle inequality and the LLN,
\begin{equation}
    \left\|\hat{Q} - Q\right\|_{\text{op}} = O_p\left(n^{\frac{-s}{2s+d} \wedge -\frac{1}{2}}\right).
\end{equation}
Finally we have
\begin{align}
    \left\|\hat{Q}^{-1} - Q^{-1}\right\|_{\text{op}} &= \left\|Q^{-1}\left(Q - \hat{Q}\right)\hat{Q}^{-1}\right\|_{\text{op}} \\
    &\leq \left\|Q^{-1}\right\|_{\text{op}}\left\|Q - \hat{Q}\right\|_{\text{op}}\left\|\hat{Q}^{-1}\right\|_{\text{op}}\\
    &= O_p\left(n^{\frac{-s}{2s+d} \wedge -\frac{1}{2}}\right).
\end{align}
\qed{}

% \begin{restatable}{lemma}{} \label{lemma: newey lemma}
%     We have
%     \[\frac{1}{\sqrt{n}}\sum_{i = 1}^n(\hat{w}_i - w_i)X_iu_i = -\frac{1}{\sqrt{n}}\sum_{i = 1}^n\left[\frac{p^*(x_i;\theta_0, \epsilon)(\mathbb{E}[y_i|x_i] - X_i^T\beta^*_{x_0, \epsilon})}{p_0(x_i)}\right] + o_p(1).\]
% \end{restatable}
% \noindent \textit{Proof.} Expanding the left hand side yields,
% \begin{align*}
%     \frac{1}{\sqrt{n}}\sum_{i = 1}^n(\hat{w}_i - w_i)X_iu_i &= \frac{1}{\sqrt{n}}\sum_{i = 1}^n\left[\hat{p}_0(x_i)^{-1} - p_0(x_i)^{-1}\right]a(z).
% \end{align*}
% To prove the result, it will suffice to show that the assumptions of Lemma 4.1 of \citet{Newey_Ruud_2005} are met with $a(z) = p^*(x_i;\theta_0, \epsilon) X_i\left(y_i - X_i^T\beta^*_{x_0, \epsilon}\right)$. \Cref{assumption: lce and covariate density} directly guarantees many of the assumptions, but a few remain to be verified. 
% \begin{enumerate}
%     \item We need $\mathbb{E}[a(z)\,|\,x]$ to be bounded on $C(\theta_0, \epsilon)$ (defined in \cref{eq: support}) and continuous in $x$ on a set of full Lebesgue measure. To see why this holds, observe that 
%     \[\mathbb{E}[a(z)\,|\,x] = p^*(x;\theta_0, \epsilon) X\left(\mathbb{E}[y | x] - X^T\beta^*_{x_0, \epsilon}\right).\]
%     By \cref{assumption: lce and covariate density} we know that this expression is bounded on $C(\theta_0, \epsilon)$. Continuity in $x$ almost everywhere stems from the assumption that $p^*(x; \theta_0, \epsilon)$ is continuous almost everywhere and it takes the value $0$ outside of $\mathcal{B}(0, \epsilon)$. This means that multiplication with the (discontinuous) indicator function does not get rid of continuity. 
%     \item The requirement that $a(z) = 0$ except on a compact set is trivially true. Taking that set to be $C(\theta_0, \epsilon)$ yields the requirement that $p_0(x) > 0$ for all $x \in C(\theta_0, \epsilon).$ 
%     \item Finally, we need $\mathbb{E}\left[\|a(z)\|^4\right] < \infty$. Since the $a(z)$ from \citet{Newey_Ruud_2005} differs with ours only by deletion of the indicator function, satisfaction of the condition in their case implies the same for ours. 
% \end{enumerate}
% Since all requirements of Lemma 4.1 of \citet{Newey_Ruud_2005} are satisfied, we can apply the Lemma to say
% \begin{multline}
%     \frac{1}{\sqrt{n}}\sum_{i = 1}^n(\hat{w}_i - w_i)X_iu_i = -\frac{1}{\sqrt{n}}\sum_{i = 1}^n\Bigg[\frac{p^*(x_i;\theta_0, \epsilon) X_i\left(\mathbb{E}[y_i|x_i] - X_i^T\beta^*_{x_0, \epsilon}\right)}{p_0(x_i)} \\
%     + \mathbb{E}\left[\frac{p^*(x;\theta_0, \epsilon) X\left(y - X^T\beta^*_{x_0, \epsilon}\right)}{p_0(x)}\right]\Bigg] + o_p(1).
% \end{multline}
% We'll finish the result by showing that the second term in the brackets is zero. By iterated expectation the second term in the brackets can be rewritten as 
% \begin{equation*}
%     \mathbb{E}\left[\frac{p^*(x;\theta_0, \epsilon) X\left(\mathbb{E}[y|x] - X^T\beta^*_{x_0, \epsilon}\right)}{p_0(x)}\right] = \mathbb{E}\left[\frac{p^*(x;\theta_0, \epsilon)}{p_0(x)} X\left(f(\Pi(x+x_0)) - X^T\beta^*_{x_0, \epsilon}\right)\right].
% \end{equation*}
% By the first order optimality conditions of the population optimization in \cref{eq: population estimator} this term must be zero. 

% \qed{}

\begin{restatable}[Projection norm bound]{lemma}{} \label{lemma: pnorm bound} Suppose $X\sim p_0$ and $p^*(\cdot\,;\, \theta_0)$ is some spherically symmetric distribution satisfying \cref{assumption: lce and covariate density}. Let
\[w =  \frac{p^*(x;\,\theta_0)}{p_0(x)}.\]
Then for $\epsilon$ sufficiently small we have
\begin{equation}
   \left\|\left\{\mathbb{E}_{p_{0}}[wXX^{T}]\right\}^{-1}\mathbb{E}_{p_{0}}[wX]\right\|_2 \lesssim \epsilon^{-1}.
\end{equation}
\end{restatable}

\noindent \textit{Proof}. Let $b= \Big(a^{T}\left\{\mathbb{E}_{p_{0}}[wXX^{T}]\right\}^{-1}\Big)^{T}$. We have,
\begin{align}
    \left\|\left\{\mathbb{E}_{p_0}[wXX^{T}]\right\}^{-1}\mathbb{E}_{p_0}[wX]\right\|_2 &= \sup_{\|a\|_2\neq 0} \frac{a^{T}\left\{\mathbb{E}_{p_0}[wXX^{T}]\right\}^{-1}\mathbb{E}_{p_0}[wX]}{\sqrt{a^{T}a}} \\
    &= \sup_{\|b\|_2 \neq 0} \frac{b^{T}\mathbb{E}_{p_0}[wX]}{\sqrt{b^{T} \mathbb{E}_{p_0}\big[wXX^{T}\big]\mathbb{E}_{p_0}\big[wXX^{T}\big]^{T}b}}. \label{eq: projection norm}
\end{align}
First we'll analyze the numerator. Observe that for a any $b$ we can bound it as
\begin{align}
    b^{T}\mathbb{E}_{p_{0}}[wX] &= b^{T}\,\mathbb{E}_{ p_{0}}\left[\frac{p^*(x;\theta_0, \epsilon)}{p_{0}(x)}X\right] \\
    &= b^{T}\,\mathbb{E}_{p^*(\cdot\,; \,\theta_0)}\left[X\right] \\
    &= b^{T}\int_{\mathbb{R}^d}p^*(x;\theta_0, \epsilon)X dx_1 \dots dx_d \\
    &= b^{T}\int_{\mathcal{B}(0, \epsilon)}p^*(x;\theta_0, \epsilon)X dx_1 \dots dx_d \\
    &= \int_{\mathcal{B}(0, \epsilon)}p^*(x;\theta_0, \epsilon)(b^{T}X) dx_1 \dots dx_d \\
    &\leq \epsilon \|b\|_2 \int_{\mathcal{B}(0, \epsilon)}p^*(x;\theta_0, \epsilon) dx_1 \dots dx_d \\
    &= \epsilon\|b\|_2. \label{eq: numerator bound}
\end{align}
Similarly, for a any $b$ the denominator can be bounded as
\begin{align}
    \sqrt{b^{T} \mathbb{E}_{p_0}\big[wXX^{T}\big]\mathbb{E}_{p_0}\big[wXX^{T}\big]^{T}b} &= \left\| \mathbb{E}_{p_0}\big[wXX^{T}\big]^{T}b\right\|_2 \\
    &\geq \sigma_{\min}\left(\mathbb{E}_{ p_0}\big[wXX^{T}\big]\right) \|b\|_2 \\
    &= \sigma_{\min}\left(\mathbb{E}_{p^*}\left[XX^{T}\right]\right) \|b\|_2. \label{eq: denominator bound 1}
\end{align}
Now it suffices to find a lower bound on the smallest singular value of the restricted scatter matrix. Note that since the matrix is symmetric and positive semidefinite, we have 
\begin{align*}
    \sigma_{\min}\left(\mathbb{E}_{p^*}\left[XX^{T}\right]\right) &= \lambda_{\min}\left(\mathbb{E}_{p^*}\left[\begin{pmatrix}
        1 \\ x
    \end{pmatrix}
    \begin{pmatrix}
        1 \\ x
    \end{pmatrix}^{T}\right]\right) \\
    &= \lambda_{\min}\left(\begin{bmatrix}
        1 & 0 \\
        0 & \mathbb{E}_{p^*}[xx^T]
    \end{bmatrix}\right) \\
    &= \min\{1, \nu^2\}
\end{align*}
where $\nu^2 = \mathbb{E}_{p^*}[x_i^2]$ for any $i \in [d]$. Note that $\nu^2 \asymp \epsilon^2$ as we choose $p^*$ to be spherically symmetric and supported on $\mathcal{B}(0, \epsilon)$. Plugging this, \cref{eq: denominator bound 1} and \cref{eq: numerator bound} back into \cref{eq: projection norm}, we obtain
\begin{align*}
    \left\|\left\{\mathbb{E}_{X\sim p_0}[w(X)XX^{T}]\right\}^{-1}\mathbb{E}_{X \sim p_0}[w(X)X]\right\|_2 &\leq \frac{\epsilon\|b\|_2}{\|b\|\min\{1, \nu^2\}} \\
    &= \epsilon\nu^{-2}
\end{align*}
for sufficiently small $\epsilon$. The statement follows from the fact that $\epsilon \asymp \nu$.

\qed{}


\begin{restatable}[Surrogate]{lemma}{} \label{lemma: surrogate problem} Let 
\[\tilde{\beta}^*_{x_0, \epsilon} = \arg\min_{\beta \in \mathbb{R}^{d+1}}\mathbb{E}_{ p_{0}}\Big[\frac{p^*(x;\theta_0, \epsilon)}{p_{0}(x)}\Big(g(x + x_0) - X^{T}\beta \Big)^2\Big]\]
with 
\[g(X) = f\Big(\Pi(x_0) + \nabla\Pi(x_0)^{T}(x-x_0)\Big).\]
Then for some $\alpha_0, \alpha_1 \in \mathbb{R}$ we have $\tilde{\beta}^*_{x_0, \epsilon} = \begin{bmatrix}\alpha_0  & \alpha_1\nabla \Pi(x_0) \end{bmatrix}^T$. 
\end{restatable}
\noindent \textit{Proof.} Define 
\[R_{\epsilon}(\beta; x) = \mathbb{E}_{ p_{0}}\Big[\frac{p^*(x;\theta_0, \epsilon)}{p_{0}(x)}\Big(g(x + x_0) - X^{T}\beta\Big)^2\Big].\]
Then we have,
\begin{align*}
    R_{\epsilon}(\beta; x)&= \int_{\mathbb{R}^d}\frac{p^*(x;\theta_0, \epsilon)}{p_{0}(x)}\Big(g(x+x_0) - X^{T}\beta\Big)^2p_{0}(x) dx_1 \dots dx_d \\
    &= \int_{\mathbb{R}^d}p^*(x;\theta_0, \epsilon)\Big(g(x + x_0) - X^{T}\beta \Big)^2 dx_1 \dots dx_d.
\end{align*}
Now let 
\begin{equation}
    A(\epsilon) = \int_{\mathbb{R}^d} p^*(x;\theta_0, \epsilon) dx_1 \dots dx_d.
\end{equation}
Then we have,
\begin{align*}
    R_{\epsilon}(\beta; x)&= A(\epsilon) \bigintsss_{\mathbb{R}^d}\underbrace{\frac{p^*(x;\theta_0, \epsilon)}{A(\epsilon)}}_{\tilde{p}(x;\vartheta)}\Big(g(x+x_0) - X^{T}\beta \Big)^2 dx_1 \dots dx_d \\
    &= A(\epsilon)\, \mathbb{E}_{ \tilde{p}}\left(g(x+x_0)-X^{T}\beta\right)^2.
\end{align*}
Continuing, we have
\begin{align*}
    R_{\epsilon}(\beta; X)&= A(\epsilon)\,\mathbb{E}\left[\mathbb{E}\left[\left(g(X+x_0)-X^{T}\beta\right)^2 \, \Big| 
    \, \nabla\Pi(x_0)^{T}x\right]\right]\\
    &= A(\epsilon)\,\mathbb{E}\left[\mathbb{E}\left[\left(f\Big(\Pi(x_0) + \nabla\Pi(x_0)^{T}x\Big)-X^{T}\beta\right)^2 \, \Big| 
    \, \nabla\Pi(x_0)^{T}x\right]\right].
\end{align*}
By \Cref{prop: conditional spherical symmetry} we know that $\tilde{p}$ is also a spherically symmetric distribution, and it therefore satisfies the LCE property. Now we'll apply Jensen's inequality to say,
\begin{align*}
    R_{\epsilon}(\beta; x) &\geq A( \epsilon)\,\mathbb{E}\left[\left(f\Big(\Pi(x_0) + \nabla\Pi(x_0)^{T}x\Big)-\mathbb{E}\left[X^{T}\beta \, \Big| 
    \, \nabla\Pi(x_0)^{T}x\right]\right)^2\right] \\
    &= A( \epsilon)\,\mathbb{E}\left[\left(f\Big(\Pi(x_0) + \nabla\Pi(x_0)^{T}x\Big)-\mathbb{E}\left[x^{T}\beta' \, \Big| 
    \, \nabla\Pi(x_0)^{T}x\right]- \beta_0\right)^2\right]
\end{align*}
where $\beta'$ is the last $d$ entries of $\beta$.  Since $\tilde{p}$ is spherically symmetric, it satisfies the LCE property. Thus,
\begin{align}
    R_{\epsilon}(\beta; x) &\geq A( \epsilon)\,\mathbb{E}\left[\left(f\Big(\Pi(x_0) + \nabla\Pi(x_0)^{T}x\Big)-\alpha_0' - \alpha_1\nabla\Pi(x_0)^{T}x - \beta_0\right)^2\right] \\
    &=  A( \epsilon)\,\mathbb{E}\left[\left(f\Big(\Pi(x_0) + \nabla\Pi(x_0)^{T}x\Big) - \alpha_1\nabla\Pi(x_0)^{T}x - \alpha_0\right)^2\right]. \label{eq: surrogate risk lower bound}
\end{align}
Observe that the inequality becomes an equality when $\beta = \begin{bmatrix}\alpha_0  & \alpha_1\nabla \Pi(x_0) \end{bmatrix}^T$ for some $\alpha_0, \alpha_1 \in \mathbb{R}$, which completes the proof.

\qed{}

\begin{restatable}[Proportionality constant]{rmk}{}
    The proportionality constant $\alpha_1 \in \mathbb{R}$ for \Cref{lemma: surrogate problem} may be $0$, preventing us from recovering any information about the direction of $\nabla\Pi(x_0)$. To avoid this, one must impose additional assumptions on $f$ and the distribution of the covariates. In particular, by \cref{eq: surrogate risk lower bound} we can rewrite the surrogate optimization problem as
    \begin{align}
        \alpha_0^*, \alpha_1^* &= \arg\min_{\alpha_0, \alpha_1 \in \mathbb{R}} \mathbb{E}\left[\left(f\Big(\Pi(x_0) + \nabla\Pi(x_0)^{T}x\Big) - \alpha_1 \nabla\Pi(x_0)^Tx - \alpha_0 \right)^2\right].
    \end{align}
    With some calculation one would obtain
    \begin{align}
        \alpha_1^* = \frac{\mathbb{E}\left[y\nabla\Pi(x_0)^Tx\right] - \mathbb{E}[y]\cdot \mathbb{E}\left[\nabla\Pi(x_0)^Tx\right]}{\text{Var}(\nabla\Pi(x_0)^Tx)}
    \end{align}
    where $y = f\left(\Pi(x_0) + \nabla\Pi(x_0)^Tx\right)$. Thus, $
    \alpha_1^*$ is nonzero if and only if $y \not\!\perp\!\!\!\perp \nabla\Pi(x_0)^Tx$. Equivalently, 
    \[\alpha_1^* \neq 0 \iff \mathbb{E}\left[f\left(\Pi(x_0) + \nabla\Pi(x_0)^Tx\right)\nabla\Pi(x_0)^Tx\right] \neq 0.\]
    Since the arguments in \Cref{lemma: surrogate problem} are applicable to any convex cost, one may sidestep the problem of $\alpha^*_1 = 0$ computationally by obtaining an estimate of $\nabla\Pi(x_0)$ with a different convex loss and averaging with the estimate obtained via squared loss.
    
    
\end{restatable}

% By spherical symmetry of the distribution of $x$ in the expectation above, we know that $\nabla\Pi(x_i)^Tx$ is some 1-dimensional distribution symmetric about $0$. Rewriting, we see that the condition boils down to 
% \[\mathbb{E}_z[f(c + z)z] \neq 0\]
% where $z$ is some symmetric density about $0$. Expanding
% \begin{align*}
%     \int_{\mathbb{R}}f(c+z)z p(z)dz &= \int_{-\infty}^0 f(c+z)z p(z)dz + \int_{0}^\infty f(c+z)z p(z)dz \\
%     &= \int_{0}^\infty f(c+(-z))(-z) p(-z)d(-z) + \int_{0}^\infty f(c+z)z p(z)dz\\
%     &= \int_{0}^{\infty}[f(c + z) + f(c - z)]zp(z)dz.
% \end{align*}
% By the mean value theorem $f(c + z) = f(c) + f'(\bar{z}_1(z))z$ and $f(c - z) = f(c) - f'(\bar{z}_2(z))z$ where $\bar{z}_1: [0, \infty) \rightarrow  [c, c+z]$ and $\bar{z}_2: [0, \infty) \rightarrow [c-z, c]$. Plugging back in yields
% \begin{align*}
%     \int_{\mathbb{R}}f(c+z)z p(z)dz &= \int_\mathbb{R}\left[2f(c)zp(z) + (f'(\bar{z}_1(z)) - f'(\bar{z}_2(z))zp(z)\right]dz \\
%     &= 2f(c)\int_\mathbb{R}zp(z)dz + \int_\mathbb{R}(f'(\bar{z}_1(z)) - f'(\bar{z}_2(z))zp(z)dz \\
%     &= f(c) \mathbb{E}|z| + \frac{1}{2}\mathbb{E}\left(f'(\bar{z}_1(|z|)) - f'(\bar{z}_2(|z|))\right)|z|
% \end{align*}

\begin{restatable}[Conditional spherical symmetry]{prop}{} \label{prop: conditional spherical symmetry}
    Suppose $x$ has a density $f(x)$ satisfying spherical symmetry (\Cref{def: spherical symmetry}) with mean $0$ and covariance $\sigma^2I$ for some $\sigma > 0$. Then for any $\epsilon > 0$ the distribution of the random variable $x \,| \, x \in \mathcal{B}(0, \epsilon)$ is also spherically symmetric. 
\end{restatable}
\noindent \textit{Proof.} We have that for any orthogonal $Q$ and for any Borel set $A$
\begin{align*}
    \mathbb{P}(Qx \in A \, | \, x \in \mathcal{B}(0,\epsilon) ) &= \int_{A}f(Qx\,|\,x \in \mathcal{B}(0,\epsilon))dx_1\dots dx_d \\
    &= \bigintsss_{A}\frac{f(Qx) }{\int_{\mathcal{B}(0,\epsilon)}f(x')dx'_1 \dots dx'_d}dx_1\dots dx_d \\
    &= \bigintsss_{A}\frac{f(x) }{\int_{\mathcal{B}(0,\epsilon)}f(x')dx'_1 \dots dx'_d}dx_1\dots dx_d
\end{align*}
where the last line follows from spherical symmetry of $x$. Thus,
\begin{align*}
    \mathbb{P}(Qx \in A \, | \, x \in \mathcal{B}(0,\epsilon) ) &= \bigintsss_{A}\frac{f(x) }{\int_{\mathcal{B}(0,\epsilon)}f(x')dx'_1 \dots dx'_d}dx_1\dots dx_d \\
    &= \int_{A}f(x\,|\,x \in \mathcal{B}(0,\epsilon))dx_1\dots dx_d \\
    &= \mathbb{P}(x \in A \, | \, x \in \mathcal{B}(0,\epsilon) ).
\end{align*}
\qed{}

\begin{restatable}[Approximation]{lemma}{} \label{lemma: proj approx}
    Suppose $f$ is $L_f$-Lipschitz. Then for Lebesgue almost every $x_0 \in \mathcal{B}(0,r)$, there exists an $\epsilon_0 > 0$ such that for all $\epsilon \leq \epsilon_0$ and for all $x \in \mathcal{B}(x_0, \epsilon)$
    \[\bigg|f\big(\Pi(x)\big) - f\big(\Pi(x_0) + \nabla\Pi(x_0)^{T}(x-x_0) \big)\bigg| \lesssim (M_2+rM_3)L_f\epsilon^2. \]
\end{restatable}

\noindent \textit{Proof.} By \Cref{lemma: operator norm bound} we have that for Lebesgue almost every $x_0 \in \mathcal{B}(0,r)$, there exists finite $\epsilon_0, \lambda > 0$ such that for all $\epsilon \leq \epsilon_0$ and for every $x \in \mathcal{B}(x_0, \epsilon)$
\[\|\nabla^2\Pi(x)\|_{\text{op}} \leq \underbrace{2M_2\lambda^{2}+rM_3\lambda^{3}}_{= C_{x_0, \lambda}}.\]
We can employ an expansion of the projection map as follows,
\begin{align*}
    \Pi(x) &= \Pi(x_0) + \nabla\Pi(x_0)^{T}(x-x_0) + \underbrace{\int_{0}^{1}(x-x_0)^T\nabla^2\Pi(x+t(x-x_0))(x-x_0) dt}_{= R(x)}.
\end{align*}
By the operator norm bound, we have that
\[\|R(x)\|_2 \leq \int_{0}^{1}\|x - x_0\|_2^2\|\|\nabla^2\Pi(x+t(x-x_0))\|_{\text{op}} \,dt \leq C_{x_0, \lambda}\epsilon^2.\]
It follows that 
\begin{align}
    \left|\,\Pi(x) - \left(\Pi(x_0) + \nabla\Pi(x_0)^{T}(x-x_0)\right)\right| &\leq C_{x_0, \lambda}\epsilon^2.
\end{align}
The lemma follows from the Lipschitz continuity of $f$. 

\qed{}


\begin{restatable}[Perturbation]{lemma}{} \label{lemma: perturbation}
    Let 
    \[\beta^* = \arg\min_{\beta}\mathbb{E}\big[w\big(Y - \beta^{T}X\big)^2\big]\]
    and define $\tilde{\beta}^*$ to be the solution to the same optimization with perturbed responses,
    \[\tilde{\beta}^* = \arg\min_{\beta}\mathbb{E}\big[w\big(Y - \xi(x) - \beta^{T}X\big)^2\big].\]
    If $\mathbb{E}[wXX^T] \succ 0$, then we have
    \[\beta^* - \tilde{\beta}^* = \mathbb{E}[wXX^{T}]^{-1}\mathbb{E}[wX\xi(x)].\]
\end{restatable}
\noindent \textit{Proof.} We can use the normal equations to say
\begin{align*}
    \mathbb{E}[wXX^{T}]\beta^* = \mathbb{E}[wXY]
\end{align*}
implying 
\[\beta^* = \mathbb{E}[wXX^{T}]^{-1}\mathbb{E}[wXY].\]
We can apply the same argument to obtain 
\[\tilde{\beta}^* = \mathbb{E}[wXX^{T}]^{-1}\mathbb{E}[wX(Y - \xi(x))].\]
Taking the difference, we obtain
\begin{align*}
    \beta^* - \tilde{\beta}^* &= \mathbb{E}[wXX^{T}]^{-1}\mathbb{E}[wX(Y -  Y + \xi(X) )] \\
    &= \mathbb{E}[wXX^{T}]^{-1}\mathbb{E}[wX\xi(x)].
\end{align*}
\qed{}

\begin{restatable}{lemma}{} \label{lemma: pd}
    The matrix $\mathbb{E}_{p_0}[wXX^T]$ is positive definite.
\end{restatable}
\noindent \textit{Proof.} Expanding $w$ yields
\begin{align}
    \mathbb{E}_{p_0}[wXX^T] &= \mathbb{E}_{p_0}\left[\frac{p^*(x; \theta_0, \epsilon)}{p_0(x)}XX^T\right] \\
    &= \mathbb{E}_{p^*}\left[XX^T\right] \\
    &= \mathbb{E}_{p^*}\left[XX^T\right]  \\
    &\succ 0.
\end{align}
where the last two steps follow from \cref{assumption: lce and covariate density}.

\qed{}

\begin{restatable}[Projection gradient]{prop}{} \label{prop: proj gradient}
    Let $\Pi: \mathbb{R}^d \rightarrow [0, \text{len}_{\gamma}]$ be the closest-point projection onto $\gamma$ as defined in \Cref{def: projection}. Then for every $x \notin \overline{\text{Sing}(\Gamma)}$ such that 
    \[\langle\gamma''(\Pi(x)), x - \gamma(\Pi(x))\rangle \neq 1\]
    we have 
    \[\nabla\Pi(x)  = -\frac{\gamma'(\Pi(x))}{-1 + \left\langle\gamma''(\Pi(x)), x - \gamma(\Pi(x))\right\rangle}.\]
\end{restatable}

\noindent \textit{Proof of \Cref{prop: proj gradient}.} Let $f(x,t) = \|x - \gamma(t)\|_2^2$. By non-negativity of $\|\cdot\|_2$ and monotonicity of $h(x) = x^2$ for $x \geq 0$, we have $\Pi(x) = \arg\min_t \|x - \gamma(t)\|_2 = \arg \min_t \|x - \gamma(t)\|_2^2 = \arg\min_t f(x, t)$. We have the first-order condition that 
\[ \partial_t f(x,t) = 0\]
at the optimum. Now, if we define $F(x, t) = \partial_t f(x,t)$ the first order condition amounts to the requirement that $F(x,t) = 0$. Expanding $F$, we have
\begin{align}
    F(x, t) &= \partial_t \|x - \gamma(t)\|_2^2 \\ \label{eq: ift}
    &= -2\langle x - \gamma(t)), \gamma'(t)\rangle.
\end{align}
Observe that the first order condition implies that the projection displacement $x - \gamma(t)$ lies normal to the curve at optimality. 
Now we'll apply the implicit function theorem to $F(x, t)$. Because $x$ is \textit{not} in the closure of the singular set, we have that $\Pi(x)$ is unique and $\Pi$ is differentiable in a neighborhood about $x$ - its not hard to verify that this yields continuous differentiability of $F(x,t)$ in a neighborhood about $(x, \Pi(x))$. This, coupled with the assumption that $\langle\gamma''(\Pi(x)), x - \gamma(\Pi(x))\rangle \neq 1$ allows us apply the implicit function theorem to conclude that 
\begin{align}
    \nabla\Pi(x) &= - \frac{\nabla_x F\big(x, t\big)}{\partial_t F\big(x, t\big)} \Bigg|_{t = \Pi(x)}\\
    &= - \frac{\nabla_x\left(-2\langle x - \gamma(t), \gamma'(t)\rangle\right)}{\partial_t \left(-2\langle x - \gamma(t), \gamma'(t)\rangle\right)} \Bigg|_{t = \Pi(x)}\\
    &= -\frac{\gamma'(t)}{-\|\gamma'(t)\|_2^2 + 
    \langle\gamma''(t), x - \gamma(t)\rangle}\Bigg|_{t = \Pi(x)} \label{eq: projection grad} \\ 
    &= -\frac{\gamma'(\Pi(x))}{-1 + \left\langle\gamma''(\Pi(x)), x - \gamma(\Pi(x))\right\rangle}.
\end{align}
\qed{}


\begin{restatable}[Projection hessian]{prop}{} \label{prop: proj hessian}
    Let $\Pi: \mathbb{R}^d \rightarrow [0, \text{len}_{\gamma}]$ be the closest-point projection onto $\gamma$ as defined in \Cref{def: projection}. Then for every $x \notin \overline{\text{Sing}(\Gamma)}$ and 
    \[\langle\gamma''(s), x - \gamma(s)\rangle \neq 1\]
    we have
    \[\nabla^2\Pi(x) = \frac{1}{g^2(x)}\left(\gamma''(s)\gamma'(s)^T + \gamma'(s)\gamma''(s)^T\right)- \frac{\langle\gamma^{(3)}(s), x - \gamma(s)\rangle}{g^3(x)}\gamma'(s)\gamma'(s)^T.\]
    where $s = \Pi(x)$ and $g(x) = -1 + \langle\gamma''(s), x-\gamma(s)\rangle$.
\end{restatable}
\noindent \textit{Proof.} By \cref{prop: proj gradient} we know that for all $x$ such that the assumptions hold, we have
\begin{align}
    \nabla\Pi(x) &= -\frac{\gamma'(s)}{-1 + \langle\gamma''(s), x - \gamma(s)\rangle} \\
    &= -\frac{\gamma'(s)}{g(x)}.
\end{align}
Taking the gradient again yields
\begin{align}
    \nabla^2\Pi(x) &= -\frac{g(x)\nabla \gamma'(s) - \gamma'(s)\nabla g(x)^T}{g(x)^2} \\
    &= -\frac{\nabla \gamma'(s)}{g(x)} + \frac{1}{g^2(x)}\gamma'(s)\nabla g(x)^T.
\end{align}
By the chain rule, we have $\nabla\gamma'(s) = \gamma''(s)\nabla\Pi(x)^T$ and 
\begin{align}
    \nabla g(x) &= \nabla\langle\gamma''(s), x\rangle -\nabla\langle \gamma''(s), \gamma(s)\rangle \\
    &= \langle \gamma^{(3)}(s), x\rangle \nabla\Pi(x) + \gamma''(s) - \langle\gamma^{(3)}(s), \gamma(s)\rangle\nabla\Pi(x) - \underbrace{\langle\gamma''(s), \gamma'(s)\rangle}_{= 0}\nabla\Pi(x) \\
    &= \langle\gamma^{(3)}(s), x - \gamma(s)\rangle\nabla\Pi(x) + \gamma''(s).
\end{align}
Thus,
\begin{align*}
    \nabla^2\Pi(x) &= -\frac{1}{g(x)}\gamma''(s)\nabla\Pi(x)^T + \frac{\langle\gamma^{(3)}(s), x - \gamma(s)\rangle}{g^2(x)}\gamma'(s)\nabla\Pi(x)^T + \frac{1}{g^2(x)}\gamma'(s)\gamma''(s)^T \\
    &= \frac{1}{g^2(x)}\gamma''(s)\gamma'(s)^T - \frac{\langle\gamma^{(3)}(s), x - \gamma(s)\rangle}{g^3(x)}\gamma'(s)\gamma'(s)^T + \frac{1}{g^2(x)}\gamma'(s)\gamma''(s)^T \\
    &= \frac{1}{g^2(x)}\left(\gamma''(s)\gamma'(s)^T + \gamma'(s)\gamma''(s)^T\right)- \frac{\langle\gamma^{(3)}(s), x - \gamma(s)\rangle}{g^3(x)}\gamma'(s)\gamma'(s)^T.
\end{align*}
\qed{}

\begin{restatable}[Operator norm bound]{lemma}{} \label{lemma: operator norm bound}
    Suppose $x \in \mathcal{B}(0,r) \cap R$ with $R$ defined in \cref{eq: good set}. Then there exist finite $\epsilon_0, \lambda > 0$ such that for every $\epsilon \leq \epsilon_0$ and every $y\in \mathcal{B}(x,\epsilon)$ 
    \[\|\nabla^2\Pi(y)\|_{\text{op}} \leq 2M_2\lambda^{2}+rM_3\lambda^{3}.\]
\end{restatable}

\noindent \textit{Proof.} By \Cref{thm: differentiability of pi} we have that for every $x \in R$, 
\[\nabla^2\Pi(x) = \frac{1}{g^2(x)}\left(\gamma''(s)\gamma'(s)^T + \gamma'(s)\gamma''(s)^T\right)- \frac{\langle\gamma^{(3)}(s), x - \gamma(s)\rangle}{g^3(x)}\gamma'(s)\gamma'(s)^T.\]
where $s = \Pi(x)$ and $g(x) = -1 + \langle\gamma''(s), x-\gamma(s)\rangle$. Since $R$ is an open set, there exists an $\epsilon_0 > 0$ such that this expression holds for all $y \in \mathcal{B}(x, \epsilon)$ for $\epsilon \leq \epsilon_0$.  Bounding the operator norm yields
\begin{align*}
    \|\nabla^2\Pi(x)\|_{\text{op}} &= \max_{v\in \mathbb{S}^{d-1}} \|\nabla^2\Pi(x)v\|_2 \\
    &\leq \max_{v\in \mathbb{S}^{d-1}}\left(\frac{\|\gamma''(s)\gamma'(s)^Tv + \gamma'(s)\gamma''(s)^Tv\|_2}{|g^2(x)|} + \frac{\langle\gamma^{(3)}(s), x - \gamma(s)\rangle}{|g^3(x)|} \|\gamma'(s)\gamma'(s)^Tv\|_2\right) \\
    &\leq \max_{v\in \mathbb{S}^{d-1}}\left(\frac{\|\gamma''(s)\gamma'(s)^Tv\|_2 + \|\gamma'(s)\gamma''(s)^Tv\|_2}{|g^2(x)|} + \frac{\langle\gamma^{(3)}(s), x - \gamma(s)\rangle}{|g^3(x)|} \|\gamma'(s)\gamma'(s)^Tv\|_2\right) \\
    &\leq \frac{2M_2}{|g^2(x)|} + \frac{\langle\gamma^{(3)}(s), x - \gamma(s)\rangle}{|g^3(x)|} \\
    &\leq \frac{2M_2}{|g^2(x)|} + \frac{r M_3}{|g^3(x)|}.
\end{align*}
where the final line follows from the assumption that $\Gamma = \text{Im}(\gamma) \subset \mathcal{B}(0, r)$ and $x \in \mathcal{B}(0, r)$. Now we will establish that $1/|g(x)|$ is continuous everywhere in $\mathcal{B}(0,r) \cap R$. This property is a consequence of 
\begin{enumerate}
    \item the of $|g(x)|$ in $\mathcal{B}(0,r) \cap R$, which follows from the definitions of $g$ and $R$
    \item the fact that $|g(x)| \neq 0$ for all $x \in R$ -- this too follows from the definitions of $g$ and $R$.
\end{enumerate}

From these two properties, we can deduce that $1/|g(x)|$ is continuous everywhere in $\mathcal{B}(0,r) \cap R$. The continuity of $1/|g(x)|^k$ for $k\geq 1$ follows. We will leverage what has been shown to say that for every $x$ in $\mathcal{B}(0,r) \cap R$, $\alpha = 1/|g(x)| <\infty$. Furthermore, for the $\epsilon$ defined earlier, there exists a $\delta>0$ such that for all $y \in \mathcal{B}(x, \epsilon)$, $1/|g(y)| \leq \alpha+\delta = \lambda.$ The statement follows.

\qed{}

\begin{restatable}[IFT degeneracy]{lemma}{} \label{lemma: IFT degen}
    Let 
    \begin{equation} 
    B= \{x \in \mathbb{R}^d : 1 - \langle\gamma''(t), x-\gamma(t)\rangle = 0 \text{ for all }t \in \arg\inf_t \|x-\gamma(t)\|_2\} \label{eq: ift degeneracy set}
    \end{equation}
    and let $\mu$ denote the $d$-dimensional Lebesgue measure on $\mathbb{R}^d$. Then $\mu(B) = 0.$
\end{restatable}
\noindent \textit{Proof of \Cref{lemma: IFT degen}.} Let $N\gamma = \{(p, v) : p \in \gamma([0, \text{len}_{\gamma}]), v \in (T_p \gamma)^{\bot}\}$ denote the normal bundle of the curve $\gamma$, where $T_p\gamma$ denotes the tangent space of $\gamma$ at point $p$. Observe that $N\gamma$ forms a smooth $d$-dimensional manifold - explicitly, one has the canonical local chart $(t,w) \rightarrow (\gamma(t), w_1n_1(t) + \dots+w_{d-1}n_{d-1}(t))$ where $\{n_i(t)\}_{i\in [d-1]}$ form an orthonormal basis of $N_p\gamma$. 

Now let $B'= \{(p,v) \in N\gamma : \langle v,  \gamma''(\gamma^{-1}(p))\rangle = 1\}$, and let $\Phi:N\gamma \rightarrow \mathbb{R}$ be the function
\[\Phi(p,v) = \langle v, \gamma''(\gamma^{-1}(p))\rangle - 1.\]
Observe that for all $x' \in B$, for all $t' \in \arg\min_t \|x-\gamma(t)\|$ we have \[\langle \underbrace{x' - \gamma(t')}_{v'}, \gamma''(t')\rangle = 1\] where $v' \in N_{\gamma(t')}\gamma$. Thus $x' \in \exp(B')$ where $\exp(p,v) = p + v$. It follows that $B\subset \exp(B')$.

We will now show that $0$ is a regular value of $\Phi$, allowing us to conclude that the level set $\Phi^{-1}(0) = B$ is an embedded submanifold of $N\gamma$ of codimension 1. To show that $0$ is a regular value of $\Phi$, observe that $\Phi$ can be expressed in local coordinates as
\[\Phi(t,w) = \sum_{i = 1}^{d-1}w_i\langle n_i(t), \gamma''(t)\rangle - 1.\]
Taking the partial derivative with respect to $w_i$ we have
\[\partial_{w_i}\Phi(t,w) = \langle n_i(t), \gamma''(t)\rangle\]
To show the regularity of $\Phi^{-1}(0)$, we need to establish that for all $(t,w)$ such that $\Phi(t,w) = 0$, the gradient is nonzero. Observe that for all such $(t,w)$, $\langle v(t,w), \gamma''(t)\rangle = 1$, where $v(t,w) = \sum_{i = 1}^{d-1}w_in_{i}(t)$ -- therefore $\gamma''(t) \neq 0$. Now we need to guarantee that $\gamma''(t)$ has a component in $N_{\gamma(t)}\gamma$. We will show this by contradiction: suppose $\gamma''(t) \in T_{\gamma(t)}\gamma$. Then $\langle\gamma''(t), v\rangle = 0$ for all $v$ normal to $\gamma$ at $t$, which contradicts $\langle \gamma''(t), v(t,w) \rangle = 1$.

Since $\gamma''(t)$ necessarily has a some component normal to the curve, we know that $\partial_{w_i}\Phi(t,w)$ is nonzero for some $i$. This renders the differential of $\Phi$ surjective for all $(t,w)$ such that $\Phi(t,w) = 0$. This means one can apply the regular level set theorem to conclude that $B'$ forms a $d-1$ dimensional submanifold of the normal bundle of $\gamma$ (\citet{lee2003smooth}, Corollary 5.14). Now since $\exp\big|_{B'}$ is a smooth map from $B'$ to $\mathbb{R}^d$, $\exp(B')$ must have $d$-dimensional Lebesgue measure $0$ (\citet{lee2003smooth} Corollary 6.11). Since $B \subset \exp(B')$, by monotonicity of the Lebesgue measure $B$ must have measure $0$ as well.

\qed{}


\begin{restatable}[IFT degeneracy is closed]{lemma}{} \label{lemma: IFT degen closed}
    The set 
    \begin{equation}
    B = \{x \in \mathbb{R}^d : 1 - \langle\gamma''(t), x-\gamma(t)\rangle = 0 \text{ for all }t \in \arg\inf_t\|x-\gamma(t)\|_2\}
    \end{equation}
    is a closed subset of $\mathbb{R}^d$. 
\end{restatable}

\noindent \textit{Proof.} Define the set
\begin{equation}
S= \{(x,t) \in \mathbb{R}^d\times I : t \in \arg\inf_{t\in I}\|x-\gamma(t)\|_2, \, 1 - \langle\gamma''(t), x - \gamma(t)\rangle = 0\} \label{eq: S}
\end{equation}
which is trivially a subset of $\mathbb{R}^{d+1}$. If one fixes $s \in I$, the map $(x,t) \mapsto \|x-\gamma(t)\|_2 - \|x-\gamma(s)\|_2$ is continuous. Since sublevel sets of continuous functions are closed, the set
\begin{equation}
S_s = \{(x,t) : \|x-\gamma(t)\|_2 - \|x-\gamma(s)\|_2 \leq 0\}
\end{equation}
is a closed subset of $\mathbb{R}^{d+1}$. Observe that
\begin{align*}
    \bigcap_{s \in I} S_s &= \bigcap_{s \in I} \Big\{(x,t) : \|x-\gamma(t)\|_2 - \|x-\gamma(s)\|_2 \leq 0\Big \} \\
    &= \Big\{(x,t) : t\in \arg\inf_{t\in I}\|x-\gamma(t)\|_2 \Big \} \\
    &= S'.
\end{align*}
To see why consider the following. The element $(x,t)$ is in $S'$ if and only if $\|x - \gamma(t)\|_2  \leq \|x - \gamma(s)\|_2$ for all $s \in I$. Since the image of $\gamma$ is closed, this occurs if and only if $t \in \arg\inf_{t \in I} \|x - \gamma(t)\|_2$. Having established the closure of $S_s$, the closure of $S'$ follows from the fact that uncountable intersections of closed sets are still closed. 

To establish the closure of the set $S$ defined in \cref{eq: S}, consider the map $(x,t) \mapsto 1 - \langle \gamma''(t), x - \gamma(t)\rangle$. Since $\gamma \in C^3(I)$ this is a continuous map. It follows that the level set $\{(x,t) : 1 - \langle \gamma''(t), x - \gamma(t)\rangle = 0\}$ is closed. Taking the intersection with $S'$ implies the closure of the set $S$. 

Now we will use this to show that $B$ is a closed subset of $\mathbb{R}^d$. To see this, observe that $B$ can be rewritten as $B = \{x : (x,t) \in S\}.$ We will show that $B$ is closed by showing that it contains all of its limit points via its connection to the closed set $S$. Let $\tilde{x}$ be a limit point of $B$. This means there exists a sequence $x_n \in B$ such that $x_n\rightarrow \tilde{x}$. By the connection to $S$ established above, we know that there are associated values $t_n$ such that 
\[t_n \in \arg\inf_{t \in I}\|x_n - \gamma(t_n)\|_2 \quad \text{and} \quad 1 - \langle\gamma''(t_n), x_n - \gamma(t_n)\rangle = 0.\]
Thus, we have a sequence $(x_n, t_n) \in S$ such that $x_n \rightarrow x$. Since $I$, the domain of $\gamma$, is compact, we know that there must exist a subsequence $t_{n_k}$ such that $t_{n_k}\rightarrow t$ for some $t \in I$. Since $x_n$ converges to $x$, we know that the subsequence $x_{n_k}$ also converges to $x$. Thus, we have a sequence $(x_{n_k}, t_{n_k}) \rightarrow (x,t)$, which must be in $S$ due to closure. It follows that $x$ must also be in $B$, and $x_{n_k} \rightarrow x$. Since every subsequence of a convergent sequence converges to the same value, we know that the limit point of $x_{n}$ (which we defined earlier as $\tilde{x}$) is equal to $x$, which is in $B$. Thus, $B$ necessarily contains all of its limit points, and it is therefore closed.  

\qed{}


\begin{restatable}[Closure of the singular set has zero Lebesgue measure]{lemma}{} \label{lemma: singular set measure}
    The closure of the set 
    \begin{equation}\text{Sing}(\Gamma) = \{x \in \mathbb{R}^d: \Pi(x) \text{ fails to be differentiable}\} \label{eq: singular set}
    \end{equation}
    has $d$-dimensional Lebesgue measure $0$. 
\end{restatable}
\noindent \textit{Proof.} Let $d(x, \Gamma)$ be the distance from $x$ to the set $\Gamma$. We will first show that the closure of the set where $d^2(x,\Gamma)$ fails to be differentiable has measure zero. Call this set $\text{Sing}'(\Gamma)$. To show this, observe that the image of the curve $\Gamma$ forms a 1-dimensional $C^3$ submanifold of $\mathbb{R}^d$ with boundary
\[\partial\Gamma = \{\gamma(0), \gamma(\text{len}_{\gamma})\}.\]
Since the boundary consists of a union of two disjoint points, it forms a smooth $0$-dimensional submanifold of $\mathbb{R^d}$. This allows us to apply Remark 4.2 and Corollary 4.12 of \citet{mantegazza2002hamilton} to conclude that the Hausdorff dimension of $\overline{\text{Sing}'(\Gamma)}$ is at most $d-1$. This implies that its $d$-dimensional Hausdorff measure is $0$, which in turn implies that its $d$-dimensional Lebesgue outer measure is $0$. Since the set $\overline{\text{Sing}'(\Gamma)}$ is closed, it is necessarily Borel. It follows that the set is Lebesgue measurable with measure $0$. 

Now we will show that $\Pi(x)$ failing to be differentiable implies $d^2(x, \Gamma)$ fails to be differentiable. To see this observe that
\begin{align}
    d^2(x, \Gamma) &= \inf_{p \in \Gamma}\|x - p\|_2^2 \\ \label{eq: distance to Gamma}
    &= \|x - \gamma(t')\|_2^2
\end{align}
for any $t' \in \Pi(x)$. If $\Pi(x)$ is a singleton, then $d^2(x,\Gamma) = \|x - \gamma(\Pi(x))\|_2^2$ and we have that $d^2(x,\Gamma)$ is not differentiable at $x$ if $\Pi$ is not differentiable at $x$. Now suppose $\Pi(x)$ is not a singleton. We will show that this implies $d^2(x,\Gamma)$ is not differentiable by contradiction. When $\Pi(x)$ is not a singleton then $x$ has at least two distinct nearest points -- pick any two and call them $p$ and $q$. If $d^2(x,\Gamma)$ were differentiable at $x$ then its gradient would be unique and equal to $(x-p')$ where $p'$ realizes the infimum above. Since $p$ and $q$ both realize the infimum, the uniqueness of the gradient of $d^2(x, \Gamma)$ implies that $(x - p) = (x-q)$, which is a contradiction.

This establishes that if $\Pi$ fails to be differentiable at $x$, then $d^2(x, \Gamma)$ must also fail to be differentiable at $x$. This implies $\text{Sing}(\Gamma) \subseteq \text{Sing}'(\Gamma)$, which in turn implies the closure of the former is contained in the closure of the latter. Earlier we established that $\overline{\text{Sing}'(\Gamma)}$ has measure zero -- the Lemma statement therefore follows from monotonicity of the Lebesgue measure.

\qed{}



% \begin{restatable}[Medial axis is closed]{lemma}{}
%     The medial axis of the curve $\gamma$ is closed. 
% \end{restatable}
% \noindent \textit{Proof.} The medial axis $M(\Gamma)$ of $\Gamma = \text{Im}(\gamma)$ can be expressed as
% \[M(\Gamma) = \{x \in \mathbb{R}^d : |\{y \in \Gamma : \|x-y\|_2 = d(x,\Gamma)\}| > 1\}.\]
% To show that $M(\Gamma)$ is closed, we will show that it contains all of its limit points. Consider any sequence $x_n \in M(\Gamma)$ and suppose $x_n \rightarrow x$. We will show that $x$ is necessarily in $M(\Gamma)$. To do so, observe that $\Gamma \times \Gamma$ is compact as $\Gamma$ is compact. Since each $x_n$ is in $M(\Gamma)$, we know it has at least two closest points; we'll pick any two and call them $q_n$ and $p_n$ - this defines two separate sequences $\{q_n\}_{n \in \mathbb{N}}$ and $\{p_n\}_{n \in \mathbb{N}}$, but it also defines the joint sequence $\{(q_n, p_n)\}_{n\in \mathbb{N}} \subset \Gamma\times \Gamma$. By the compactness of this product space, there must exist a convergent subsequence $(q_{n_k}, p_{n_k}) \rightarrow (q,p) \in \Gamma \times \Gamma$. This, in turn, induces a subsequence $x_{n_k}$, which necessarily converges to $x$ since we assume that $x_n$ is a convergent sequence. Now we'll consider two cases. 
% \begin{enumerate}
%     \item \textbf{Case 1:} $(p \neq q)$. Observe that we have the equalities
%     \[d(x_{n_k}, \Gamma) = \|x_{n_k} - p_{n_k}\|_2 = \|x_{n_k} - q_{n_k}\|_2\]
%     by definition. By continuity, this implies that as $k \rightarrow \infty$,
%     \[d(x,\Gamma) = \|x - p\|_2 = \|x - q\|_2\]
%     with $p \neq q$. This directly implies that $x \in M(\Gamma)$. 
%     \item \textbf{Case 2:} $(p = q)$. We will show that, under the hypothesis of positive reach, this cannot happen. Since $x_{n_k} \in M(\Gamma)$, we know that $d(x_{n_k}, \Gamma) \geq \tau_{\gamma}$ for all $k$. This implies that $d(x, \Gamma) \geq \tau_{\gamma}$ as well. 
%     \begin{enumerate}
%         \item $[]$
%     \end{enumerate}
% \end{enumerate}


% Let $f:\mathbb{R}^d\times I\rightarrow \mathbb{R}$ be defined as $f(x,t) = \|x-\gamma(t)\|_2^2$. Then the medial axis $M$ can be expressed as
% \[M = \{x \in \mathbb{R}^d : |\arg\min_{t}f(x,t)| > 1\}.\]
% To show that $M$ is closed, consider the set 
% \[P = \{(x,t,s) \in \mathbb{R}^d \times I \times I: t\neq s, f(x,t) = f(x,s) \leq f(x,u) \, \forall \, u \in I\}.\]
% Observe that $(x,t,s) \in P$ if and only if $t$ and $s$ are distinct parameters of minimal distance points to $x$. Now consider the set
% \[U_{1,u} = \{(x,t,s): f(x,t) - f(x,u) \leq 0\}.\]
% For a fixed $u$, the map $(x,t,s) \mapsto f(x,t) - f(x,u)$ is continuous since $\gamma \in C^3(I)$, rendering the sublevel set closed. Now observe that
% \begin{align}
% U_1 &= \bigcap_{u \in I}U_{1,u} \\
% &= \left\{(x,t,s) : f(x,t) \leq f(x,u) \, \forall \, u \in I\right\}
% \end{align}
% is closed, since uncountable intersections of closed sets are closed. By a similar argument, one can show that the set 
% \[U_2 = \{(x,t,s): f(x,s) \leq f(x,u) \, \forall \, u \in I \}\]
% is also closed. Finally, define 
% \[U_3 = \{(x,t,s): f(x,t) = f(x,s)\}\]
% which is also closed, as it consists of the zero level set of a continuous function. 

% \qed{}

% \begin{restatable}[Continuous differentiability of $\Pi$]{lemma}{}
%     For any $x$ such that $\Pi(x)$ is differentiable in a neighborhood about $x$, $\Pi(x)$ must also be \textit{continuously} differentiable on that neighborhood.
% \end{restatable}
% \noindent \textit{Proof.} Let $d^2(x,\Gamma)$ be defined as in \cref{eq: distance to Gamma}. Under the hypothesis that $\Pi(x)$ is differentiable in a neighborhood $U$ about $x$, we know that it must also be unique and well defined on $U$. It follows that
% \begin{align}
%     d^2(x,\Gamma) = \|x - \gamma(\Pi(x))\|_2^2 \text{ for all }x \in U.
% \end{align}
% The hypothesis of the lemma coupled with the fact that $\gamma \in C^3(I)$ implies that $d^2(x,\Gamma)$ is differentiable on $U$. 


\begin{restatable}[Covariance of $\text{Unif}(\mathcal{B}(0,\epsilon))$]{prop}{} \label{prop: spherical cov}
    Suppose we have $x$ distributed according to a uniform probability measure over $\mathcal{B}(0,\epsilon)$ in $\mathbb{R}^d$. Then the covariance matrix of $x$ is $\frac{\epsilon^2}{d+2}I_d$. 
\end{restatable}
\noindent \textit{Proof.} By symmetry, the covariance matrix must be a multiple of $I_d$. To find the constant, we need to understand the distribution of $\mathbb{E}\|x\|_2$ - thus we will compute the radial pdf. Note that the cdf of $\|x\|_2$ can be written as
\begin{align*}
    \mathbb{P}(\|x\|_2 \leq r) &= \frac{\mu(\mathcal{B}(0,r))}{\mu(\mathcal{B}(0,\epsilon))} \\
    &= \left(\frac{r}{\epsilon }\right)^d
\end{align*}
where $\mu$ denotes the $d$-dimensional Lebesgue measure on $\mathbb{R}^d$. The radial pdf can be obtained by differentiating the cdf with respect to $r$,
\begin{align*}
    f_r(r) &= \frac{d}{dr}\mathbb{P}(\|x\|_2\leq r) \\
    &= \frac{dr^{d-1}}{\epsilon^d}.
\end{align*}
Now we can compute the covariance scaling factor as 
\begin{align*}
    \mathbb{E}[x_i^2] &= \frac{1}{d}\mathbb{E}\|x\|_2^2 \\
    &=\frac{1}{\epsilon^d}\int_{0}^{\epsilon}r^2r^{d-1}dr \\
    &= \frac{1}{\epsilon^d}\frac{1}{d+2}r^{d+2} \Bigg|^{r=\epsilon}_{r=0} \\
    &= \frac{\epsilon^2}{d+2}.
\end{align*}
\qed{}


\bibliography{main}

\end{document}
